\documentclass[10pt,a5paper,twoside]{article}
\usepackage[a5paper,inner=18mm,outer=15mm,top=18mm,bottom=20mm,headsep=5mm,footskip=9mm]{geometry}
\usepackage[spanish,es-nodecimaldot]{babel}
\usepackage{fontspec}
\defaultfontfeatures{Ligatures=TeX}
\setmainfont[Path=/usr/share/fonts/truetype/ebgaramond/,UprightFont=EBGaramond12-Regular.ttf,BoldFont=EBGaramond12-Bold.ttf,ItalicFont=EBGaramond12-Italic.ttf]{EB Garamond}
\setsansfont{Liberation Sans}
\usepackage{mathtools}
\usepackage{unicode-math}
\setmathfont[Path=/usr/share/texmf/fonts/opentype/public/lm-math/]{latinmodern-math.otf}
\usepackage{microtype}
\usepackage{booktabs,array,tabularx,longtable}
\usepackage{enumitem}
\usepackage{titlesec}
\usepackage{fancyhdr}
\usepackage{hyperref}
\usepackage{xurl}
\usepackage{bookmark}
\usepackage{ragged2e}
\usepackage{xcolor}
\usepackage{needspace}
\usepackage{etoolbox}
\definecolor{RBgray}{gray}{0.32}
\hypersetup{colorlinks=true,linkcolor=black,urlcolor=RBgray,citecolor=black,pdftitle={Reproducción social, secuenciación monetaria y necesidad de adelanto — v1.02},pdfauthor={Raúl Alberto Rivera Boullosa},pdflang={es}}
\setlength{\parindent}{1.15em}
\setlength{\parskip}{0pt}
\linespread{1.035}
\emergencystretch=2.5em
\raggedbottom
\clubpenalty=10000\widowpenalty=10000\displaywidowpenalty=10000
\setcounter{secnumdepth}{2}
\renewcommand{\thesection}{\Roman{section}}
\renewcommand{\thesubsection}{\alph{subsection}}
\titleformat{\section}{\normalfont\Large}{\thesection.}{.55em}{}
\titleformat{name=\section,numberless}{\normalfont\Large}{}{0pt}{}
\titleformat{\subsection}{\normalfont\large\itshape}{\thesubsection.}{.5em}{}
\titleformat{name=\subsection,numberless}{\normalfont\large\itshape}{}{0pt}{}
\titleformat{\subsubsection}{\normalfont\normalsize\itshape}{}{0pt}{}
\titlespacing*{\section}{0pt}{2.2ex plus .4ex}{1.1ex}
\titlespacing*{\subsection}{0pt}{1.7ex plus .3ex}{.7ex}
\titlespacing*{\subsubsection}{0pt}{1.3ex plus .2ex}{.5ex}
\pretocmd{\section}{\Needspace{6\baselineskip}}{}{}
\pretocmd{\subsection}{\Needspace{4\baselineskip}}{}{}
\usepackage{emptypage}
\pagestyle{fancy}\fancyhf{}
\fancyhead[LE]{\footnotesize\itshape RB · Ensayos 2026 · 02}
\fancyhead[RO]{\footnotesize\itshape\nouppercase{\rightmark}}
\fancyfoot[LE,RO]{\footnotesize\thepage}
\renewcommand{\headrulewidth}{.3pt}
\renewcommand{\footrulewidth}{0pt}
\makeatletter
\renewcommand\footnotesize{\@setfontsize\footnotesize{8pt}{9.4pt}}
\makeatother
\newenvironment{refblock}{%
  \begingroup\fontsize{8.8pt}{10.4pt}\selectfont\RaggedRight\setlength{\emergencystretch}{1em}
  \setlist[itemize]{label={},leftmargin=1.2em,labelsep=0pt,itemindent=-1.2em,itemsep=.25ex,topsep=.35ex,parsep=0pt,partopsep=0pt}%
}{\endgroup}
\urlstyle{same}
\setcounter{tocdepth}{2}
\addto\captionsspanish{\renewcommand{\contentsname}{Contenido}}
\usepackage{tocloft}
\setlength{\cftbeforesecskip}{.45ex}
\setlength{\cftbeforesubsecskip}{0pt}
\renewcommand{\cftsecfont}{\normalfont}
\renewcommand{\cftsecpagefont}{\normalfont}
\renewcommand{\cftsubsecfont}{\small\itshape}
\renewcommand{\cftsubsecpagefont}{\small}
\setlength{\cftsubsecindent}{1.4em}
\setlength{\cftsubsecnumwidth}{2.2em}
\newcommand{\RBblankpage}{\clearpage\thispagestyle{empty}\null\clearpage}
\newcommand{\RBfronttitle}[1]{\phantomsection\addcontentsline{toc}{section}{#1}}
\providecommand{\tightlist}{\setlength{\itemsep}{0pt}\setlength{\parskip}{0pt}}
\author{}\date{}
\begin{document}

% --- Preliminares del cuaderno ---
\begin{titlepage}\thispagestyle{empty}
\vspace*{18mm}
{\centering
{\small\scshape RB · Ensayos 2026 · 02\par}
\vspace{16mm}
{\fontsize{21}{23}\selectfont Reproducción social, secuenciación\\[2mm] monetaria\\[2mm] y necesidad de adelanto\par}
\vspace{11mm}
{\large Raúl Alberto Rivera Boullosa\par}
\vfill
{\small Facultad de Economía\\Universidad Nacional Autónoma de México\par}
\vspace{10mm}
{\small\scshape Ciudad de México · 2026\par}}
\end{titlepage}

\thispagestyle{empty}\null\clearpage

% Portadilla
\thispagestyle{empty}
\vspace*{42mm}
\begin{center}
{\small\scshape RB · Ensayos 2026 · 02\par}
\vspace{9mm}
{\Large Reproducción social, secuenciación monetaria\\y necesidad de adelanto\par}
\end{center}
\clearpage

% Verso de portadilla / datos editoriales
\thispagestyle{empty}
\vspace*{\fill}
{\footnotesize\RaggedRight
Raúl Alberto Rivera Boullosa\par
\textit{Reproducción social, secuenciación monetaria y necesidad de adelanto}.\par
RB · Ensayos 2026 · 02.\par
Facultad de Economía, Universidad Nacional Autónoma de México.\par
Ciudad de México, octubre de 2026.\par
v1.02 · edición revisada.\par
}
\clearpage

% Portada interior
\thispagestyle{empty}
\vspace*{25mm}
\begin{center}
{\small\scshape RB · Ensayos 2026 · 02\par}
\vspace{13mm}
{\fontsize{19}{21}\selectfont Reproducción social, secuenciación monetaria\\[1mm] y necesidad de adelanto\par}
\vspace{10mm}
{\large Raúl Alberto Rivera Boullosa\par}
\vfill
{\small Facultad de Economía · UNAM\par}
{\small Ciudad de México · 2026\par}
\end{center}
\clearpage

\tableofcontents
\thispagestyle{empty}
\clearpage

\section*{Resumen general}\label{resumen-general}\markright{Resumen general}\RBfronttitle{Resumen general}

Los esquemas de reproducción establecen relaciones de reposición,
distribución y continuidad, pero no determinan por sí mismos la secuencia
monetaria mediante la cual esas relaciones se realizan. Este volumen estudia
qué información adicional hace falta para pasar de una estructura
reproductiva a su ejecución monetaria y qué resultados cambian cuando cambia
el orden de esa ejecución.

La primera diferencia aparece entre el resultado terminal y el camino que
conduce a él: una misma red de obligaciones puede terminar en los mismos
saldos y recorrer, sin embargo, trayectorias distributivas distintas, cuya
profundidad depende del orden de los pagos; conocerla todavía no determina la
disponibilidad necesaria, pues hace falta especificar qué posiciones pueden
consolidarse. Una vez fijada esa
arquitectura, el requerimiento resultante tampoco equivale necesariamente a
una exposición efectiva, pues ésta depende de los recursos heredados con que
comienza el período. Y una exposición dada puede ser cubierta mediante
contratos e instrumentos distintos, capaces a su vez de dejar estados
financieros terminales diferentes.

Los artículos I--V reconstruyen sucesivamente estas mediaciones: de la
estructura reproductiva a las obligaciones y su ejecución; de ésta al
requerimiento bajo distintas arquitecturas de consolidación; y luego a la
exposición, la cobertura, la implementación y el estado financiero heredable.
El artículo VI responde a un problema que reaparece a lo largo de esa
construcción: cuándo la información conservada por una representación basta
para determinar otro observable. El artículo VII vuelve finalmente a la
reproducción y plantea el problema que queda abierto: bajo qué condiciones las
diferencias producidas por la ejecución intraperiódica sobreviven a la
frontera del período y modifican el estado reproductivo siguiente.

El resultado es una teoría por capas de determinación condicionada: cada capa
responde una pregunta precisa y el paso a la siguiente requiere información
adicional; así puede establecerse qué determina la estructura reproductiva,
qué depende de la secuencia de ejecución, qué cambia con las instituciones de
consolidación y financiación y qué permanece todavía indeterminado al cerrar
el período.

\begin{center}\rule{0.5\linewidth}{0.5pt}\end{center}

\cleardoublepage\section*{Introducción general}\label{introducciuxf3n-general}\markright{Introducción general}\RBfronttitle{Introducción general}\thispagestyle{plain}

Los esquemas de reproducción describen relaciones de reposición,
distribución y continuidad. Pueden establecer qué magnitudes deben volver a
producirse, qué sectores entregan y reciben productos, qué ingresos se
distribuyen y qué relaciones agregadas son compatibles con la continuidad del
sistema (Marx, 1976 {[}1885{]}; Trigg, 2006; Gehrke, 2018). Para ejecutar
monetariamente esas relaciones hace falta, sin embargo, información
adicional. El problema económico que organiza el volumen puede formularse así: ¿cómo
puede realizarse monetariamente una reproducción social cuando los pagos y cobros que
la realizan no ocurren simultáneamente y los medios de pago no están necesariamente
donde y cuando se necesitan? Las capas que siguen separan, sucesivamente, el problema
de constituir las obligaciones, el desfase entre pagos y cobros, la localización de la
disponibilidad, la exposición que queda descubierta, su financiación y la herencia que
esa ejecución deja al periodo siguiente.

La primera cuestión es determinar las obligaciones que corresponden a la
estructura reproductiva: una relación de reposición entre sectores no
establece por sí sola quién paga a quién ni bajo qué relación institucional
se constituye ese pago. La distinción tiene antecedentes en el tratamiento
del dinero como medio de pago de Marx (1975 {[}1867{]}) y en la literatura
posterior sobre circulación monetaria (Campbell, 1998). Trigg (2021)
constituye un antecedente particularmente próximo al articular reproducción,
contratos de crédito y crisis de pagos bajo circulación simple.

Supongamos ahora que las obligaciones ya han sido determinadas: todavía queda
por conocer cómo se ejecutan, pues el conjunto de pagos que debe realizarse no
fija el orden en que éstos tendrán lugar. La diferencia es irrelevante para algunas
preguntas terminales, pero resulta decisiva cuando interesa conocer las
posiciones que atraviesan los agentes antes del cierre. Dos ejecuciones pueden
contener exactamente las mismas obligaciones y terminar con los mismos
saldos, pero alcanzar déficits transitorios diferentes.

La literatura sobre sistemas de pagos ha mostrado, en contextos
institucionales distintos del presente modelo, que el orden, la
sincronización, la reutilización de entradas y el neteo pueden modificar la
liquidez intradía requerida (CPSS, 1997; McAndrews y Rajan, 2000; Bech y
Garratt, 2003; Galbiati y Soramäki, 2010). Rotemberg (2011) formula el
problema de los activos líquidos necesarios para liquidar obligaciones
interconectadas; Hayakawa et al. (2019) lo convierte en un problema explícito
de liquidación secuencial de una red dirigida, y trabajos recientes optimizan
el orden de pagos en sistemas de gran valor (McMahon et al., 2024; Astuti et
al., 2026).\footnote{El antecedente formal más próximo para el funcional bajo segregación es
Hayakawa et al. (2019; versión de conferencia 2017). Su \emph{Settlement Fund
Circulation Problem} asigna a cada vértice el fondo que debe introducirse a lo
largo de una secuencia de arcos para que cada obligación pueda liquidarse. Bajo
las convenciones de este volumen, comparar ambas definiciones da
\(p_\pi(a)=[-\min_k c_a^\pi(k)]_+=D_a^\pi\), y por tanto el fondo total de
una secuencia coincide con \(M_0(\pi;\mathcal P_{\rm sep})\). Esta
equivalencia es una identificación entre las dos construcciones, no una
reivindicación atribuida literalmente a esos autores. McMahon et al. (2024) y
Astuti et al. (2026) constituyen antecedentes recientes para la optimización
computacional del reordenamiento.} El problema puede plantearse aquí dentro de
una estructura reproductiva: una vez fijadas las obligaciones, ¿qué cambia
cuando cambia solamente su ejecución?

Sea una red cerrada de cinco agentes y doce obligaciones, con volumen nominal
\(T=42\), balanceada agente por agente. Consideremos dos ejecuciones del mismo
conjunto de obligaciones. Las necesidades distributivas máximas que producen
son

\[
D^\alpha=(4,4,4,0,6),
\qquad
D^\beta=(4,0,2,0,6).
\]

El estado terminal es el mismo; las trayectorias que conducen a él no lo son.

Este resultado plantea inmediatamente una segunda cuestión. Sumar los
déficits máximos individuales supone que cada posición debe sostenerse
separadamente. Si algunas pueden consolidarse, la disponibilidad requerida
cambia. Bajo segregación completa, las dos ejecuciones anteriores requieren
18 y 12 unidades de disponibilidad inicial. Bajo \emph{pooling} completo de
la red cerrada, ambas requieren internamente cero. La secuencia no basta, por
tanto, para determinar el requerimiento sin especificar también la
arquitectura de consolidación.

Fijemos ahora esa arquitectura: el requerimiento así obtenido parte de un
saldo inicial de referencia; si los agentes comienzan el período con recursos
heredados, sólo la parte no cubierta por esos recursos constituye exposición
efectiva. Una misma necesidad calculada puede producir así exposiciones
distintas según el estado inicial.

La exposición tampoco determina por sí sola cómo será financiada. Una
trayectoria dada puede admitir varias coberturas; una misma cobertura puede
realizarse mediante contratos, instrumentos y contrapartidas diferentes; y
esas implementaciones pueden dejar al cierre depósitos, préstamos,
vencimientos y otras posiciones financieras distintas. Para conocer qué
sobrevive al período hace falta, finalmente, una regla que transforme el
estado terminal en estado heredado.

La construcción seguida en los artículos I--V puede resumirse como

\[
\begin{aligned}
(R,I_\chi)
&\mapsto E
\mapsto \Pi_{\mathrm{adm}}(E,C)
\ni \pi \\
&\mapsto c^\pi
\mapsto D^\pi
\mapsto M_0(\pi;\mathcal P) \\
&\mapsto U
\mapsto \mathcal Y(U,\mathcal K) \\
&\mapsto Z_t^M(Q_t)
\mapsto h_{t+1}.
\end{aligned}
\]

Cada paso responde a una pregunta que el objeto anterior no resuelve por sí
solo; las flechas no tienen, por ello, un único estatuto: algunas son
definiciones; otras evalúan mapas ya fijados; otras dependen de parámetros
institucionales; otras expresan solamente compatibilidad. La secuencia
representa relaciones de determinación condicionada, no una cadena causal
automática.

Esta construcción produce además un problema común: si dos ejecuciones
distintas pueden tener el mismo \(D\) y distinto \(M_0\), o una misma
exposición admitir coberturas diferentes, ¿qué información debe conservar una
descripción para determinar el observable que interesa? El artículo VI
responde esta pregunta en términos de fibras y factorización. Su función no es
añadir otra etapa a la cadena anterior, sino establecer un criterio general de
suficiencia informacional para las distintas etapas de la construcción.

Queda entonces una última cuestión: el estado \(h_{t+1}\) conserva las
posiciones monetario-financieras que sobreviven a la ejecución, pero una
dinámica reproductiva completa exige además determinar cuándo esas diferencias
modifican stocks físicos, producción en curso, capacidad o composición
productiva. El artículo VII delimita este paso como un problema abierto. La
teoría cerrada del volumen llega hasta la herencia monetario-financiera
condicionada; la transmisión de la ejecución intraperiódica al estado
reproductivo siguiente constituye el programa que queda por desarrollar.

La función de los ejemplos seguirá el estatuto de cada resultado: las
proposiciones generales se demuestran y los ejemplos las exhiben; las
no-identificaciones se establecen mediante pares pertenecientes al mismo
dominio que conservan un valor aguas arriba y producen resultados distintos
aguas abajo. Allí donde la dinámica permanece abierta, los ejemplos sirven
para precisar qué deberá explicar una construcción posterior.

\cleardoublepage
\setcounter{subsection}{0}
\vspace*{6mm}\section*{I}\label{i.-reproducciuxf3n-obligaciones-y-secuenciaciuxf3n-monetaria}
\thispagestyle{plain}\markright{I · Reproducción, obligaciones y secuenciación monetaria}
\phantomsection\addcontentsline{toc}{section}{I · Reproducción, obligaciones y secuenciación monetaria}
{\large\itshape Reproducción, obligaciones y secuenciación monetaria\par}\medskip

\subsection*{Estructura, ejecución y necesidad distributiva}\label{estructura-ejecuciuxf3n-y-necesidad-distributiva}

\subsection{Obligaciones y ejecución}\label{obligaciones-y-ejecuciuxf3n}

Una descripción reproductiva puede fijar qué debe producirse, reponerse
y distribuirse sin fijar la secuencia monetaria mediante la cual esas
relaciones se realizan. Incluso después de constituir las obligaciones
bilaterales queda por determinar el orden en que se efectuarán.

El primer problema es distinguir lo que la reproducción exige de la forma monetaria
concreta en que esas exigencias se vuelven pagos. Llamaremos estructura reproductiva a
la información que fija qué debe producirse, reponerse o distribuirse en el dominio
considerado; todavía no contiene, por ese solo hecho, una lista de deudores y acreedores.

Sea \(\mathcal A=\{1,\ldots,N\}\) el conjunto finito de agentes y sea
\(R\) una descripción de la estructura reproductiva pertinente. La
construcción admite distintas representaciones de \(R\), que pueden
contener cantidades, stocks o relaciones materiales. Constituir
una red bilateral de pagos requiere, además, convenciones que traduzcan
esa información reproductiva en obligaciones monetarias.

La pregunta siguiente es qué relaciones sociales e institucionales convierten una
entrega o reposición necesaria para la reproducción en una obligación monetaria de
alguien frente a alguien. Para separar esa mediación de la estructura productiva, 
Sea \(\Gamma_\chi(R)\) el conjunto de configuraciones institucionales
compatibles con \(R\). Para cada \(I_\chi\in\Gamma_\chi(R)\), una regla

\[
\chi(R,I_\chi)=E
\]

produce un conjunto finito de obligaciones monetarias. \(R\) identifica
unívocamente \(E\) mediante esta mediación si y sólo si
\(\chi(R,\cdot)\) es constante sobre \(\Gamma_\chi(R)\). La compatibilidad con una estructura no selecciona una configuración
institucional. La multiplicidad de configuraciones sólo produce redes
distintas si \(\chi(R,\cdot)\) toma valores distintos sobre ellas.

El paso \(R\to E\) no consiste en una mera traducción de una
representación productiva a una monetaria. Leontief y Bródy (1993) construyen una extensión de las
relaciones insumo-producto mediante stocks y flujos monetarios; Trigg (2021),
por otra vía, vincula demandas monetarias entre productores con una matriz
formal de contratos de crédito. Estos antecedentes permiten precisar la frontera que interesa aquí. Una
representación monetaria de la interdependencia productiva no constituye
todavía una arquitectura de obligaciones ejecutables: una relación de
reposición, una celda insumo-producto o una atribución por integración
vertical no determina por sí misma quién debe pagar a quién. La constitución de \(E\) requiere información institucional adicional
contenida en \(I_\chi\): por ejemplo, unidades económicas, propiedad,
contratos y exigibilidad; por ello, una misma estructura reproductiva puede
admitir distintas constituciones de obligaciones.\footnote{Debe distinguirse la
  granularidad contractual de la granularidad de liquidación. Una obligación única de cuantía \(w\) puede liquidarse mediante varias
  transferencias sin convertirse por ello en varias obligaciones: en el
  primer caso cambia el procedimiento de ejecución sin cambiar \(E\); en
  el segundo cambia el propio conjunto de obligaciones.}

Cada obligación \(e\in E\) tiene pagador \(u_e\in\mathcal A\), receptor
\(v_e\in\mathcal A\) y magnitud \(w_e>0\). Cuando la obligación se liquida, cambia de inmediato la posición monetaria del
pagador y del receptor. Representamos ese cambio, distinto de la obligación todavía no
ejecutada, mediante su vector de transferencia.

Formalmente,

\[
g_e=w_e(\mathbf e_{v_e}-\mathbf e_{u_e})\in\mathbb R^N.
\]

Para una transferencia interna,

\[
\mathbf1^\top g_e=0.
\]

La igualdad expresa conservación de la transferencia dentro del conjunto
de agentes. La posición de cada agente durante la ejecución exige seguir la
trayectoria acumulada de las transferencias.

Constituido \(E\), queda por determinar el orden de ejecución. El problema económico
es ahora qué impide que esas obligaciones se liquiden en cualquier orden imaginable:
pueden existir precedencias productivas, contractuales, legales o documentales, además
de otras restricciones de factibilidad. Sea \(\Omega(E)\) el
universo de órdenes completos de las obligaciones distinguibles.
Escribimos

\[
C=(\prec,\mathcal F_{\mathrm{adm}}),
\]

donde \(\prec\) recoge precedencias justificadas y
\(\mathcal F_{\mathrm{adm}}\subseteq\Omega(E)\) restricciones
adicionales de factibilidad. El dominio admisible es

\[
\Pi_{\mathrm{adm}}(E,C)
=
\operatorname{Lin}(E,\prec)\cap\mathcal F_{\mathrm{adm}}.
\]

Las restricciones pueden proceder de condiciones físicas, contractuales,
legales o documentales. Se fijan antes de evaluar el resultado monetario:
la admisibilidad pertenece a la definición del dominio de ejecución.

Esta prioridad lógica es compatible con reconstrucciones temporales que
insisten en no sustituir la sucesión por una valoración simultánea
(Freeman, 1995; Kliman y McGlone, 1999; Tombazos, 2014); el dominio
\(\Pi_{\mathrm{adm}}\), en cambio, es una construcción propia del
programa.\footnote{Freeman, Kliman--McGlone y Tombazos sirven aquí como
  antecedentes de una lectura temporal y no simultánea. La regla de
  admisibilidad, el índice intraperiódico, la trayectoria \(c^\pi\) y el
  funcional \(D^\pi\) pertenecen a la construcción del programa.}

\subsection{De la ejecución a la trayectoria}\label{de-la-ejecuciuxf3n-a-la-trayectoria}

Para

\[
\pi=(\pi_1,\ldots,\pi_Q)\in\Pi_{\mathrm{adm}}(E,C),
\]

necesitamos registrar dónde queda monetariamente cada agente después de cada pago.
Para ello definimos la trayectoria acumulada

\[
c^\pi(k)
=
\sum_{j=1}^k g_{\pi_j},
\qquad
c^\pi(0)=0,
\qquad
k=0,\ldots,Q.
\]

El componente \(c^\pi_a(k)\) registra cobros menos pagos acumulados del
agente \(a\) respecto de su posición inicial. En una red interna,

\[
\mathbf1^\top c^\pi(k)=0
\qquad
\forall k.
\]

Aunque exista conservación agregada, pueden aparecer déficits individuales
cuando un agente paga antes de cobrar y otro acumula simultáneamente la
posición positiva correspondiente.

La trayectoria permite plantear una pregunta operacional: ¿qué dotación inicial
mínima permitiría al agente \(a\) recorrer esta secuencia sin que su posición acumulada
quedara por debajo de cero? La magnitud que responde esa pregunta es la siguiente.

La necesidad distributiva máxima del agente \(a\) es

\[
D^\pi_a
=
\left[-\min_{0\le k\le Q}c^\pi_a(k)\right]_+.
\]

El funcional pertenece a una familia conocida de problemas de extremos de
sumas acumuladas y, en liquidación de pagos, coincide bajo segregación con el
fondo requerido por una secuencia en Hayakawa et al. (2019). También es afín a
medidas de máxima utilización o máximo débito neto acumulado empleadas en la
literatura de liquidez intradía (BCBS, 2013; Alexandrova Kabadjova et al.,
2023; McMahon et al., 2024).\footnote{La genealogía matemática es anterior a la literatura de pagos. Abdel-Wahab y
Kameda (1978) minimizan el máximo coste acumulado de una secuencia con costes
positivos y negativos; Monma y Sidney (1979) sistematizan el método de
intercambio de pares adyacentes para problemas de secuenciación, incluyendo el
de máximo coste acumulado; Kellerer et al. (1998) estudian el \emph{Stock Size
Problem}, donde actividades consumen o suministran un recurso y se minimiza el
máximo stock requerido. Para una trayectoria escalar balanceada, la inversión
temporal y el cambio de signo transforman un máximo déficit de prefijo en un
máximo stock. Por ello, ni el extremo de prefijos, ni el intercambio adyacente,
ni la dependencia general del requerimiento respecto del orden constituyen por
sí solos la contribución específica de este volumen.}

Este antecedente respalda la interpretación económica del funcional
individual y permite distinguir tres operaciones. Reordenar pagos modifica \(\pi\) manteniendo fija la
red; compensar créditos y deudas reduce obligaciones opuestas bajo reglas
de liquidación, como en el sistema de crédito de Trigg (2021); consolidar
posiciones modifica la arquitectura \(\mathcal P\) bajo una ejecución ya
fijada. Secuencia, compensación y consolidación pueden interactuar, pero son
transformaciones distintas del problema.

Equivalentemente,

\[
D^\pi_a
=
\min\left\{
d_a\ge0:
d_a+c^\pi_a(k)\ge0\ \forall k
\right\}.
\]

\Needspace{4\baselineskip}\textbf{Demostración.} Sea \(m_a=\min_k c^\pi_a(k)\). Si \(m_a\ge0\),
\(d_a=0\) satisface todas las desigualdades y es mínimo. Si \(m_a<0\),
todo \(d_a<-m_a\) deja negativa la trayectoria en un prefijo que alcanza
\(m_a\), mientras \(d_a=-m_a\) hace no negativa toda la trayectoria. Por
tanto el mínimo es \([ -m_a ]_+\).

Así, \(D^\pi_a\) es la traslación vertical contable mínima que sostiene
la trayectoria individual. Su realización mediante dinero, crédito o
financiación exige una mediación posterior.

\subsection{Piso terminal y exceso transitorio}\label{piso-terminal-y-exceso-transitorio}

Toda ejecución completa contiene exactamente una vez cada obligación de
\(E\), de modo que

\[
c^\pi(Q)
=
\sum_{e\in E}g_e
=
\Delta(E)
\]

para cualquier \(\pi\in\Pi_{\mathrm{adm}}(E,C)\). La posición terminal
es invariante al orden.

Conviene separar dos fuentes económicamente distintas de esa necesidad: la parte que
subsiste por la posición terminal y la parte adicional que aparece sólo porque hay que
atravesar una secuencia antes de llegar al cierre.
Definimos el piso terminal

\[
B_a(E)
=
[-\Delta_a(E)]_+.
\]

Como el instante terminal pertenece al conjunto de prefijos sobre el
cual se toma el mínimo,

\[
D^\pi_a\ge B_a(E).
\]

\Needspace{4\baselineskip}\textbf{Demostración.} Para cada \(a\),

\[
\min_k c^\pi_a(k)
\le
c^\pi_a(Q)
=
\Delta_a(E).
\]

Aplicar \(x\mapsto[-x]_+\) produce la desigualdad.

Definimos entonces el exceso transitorio

\[
\Xi^\pi:=D^\pi-B(E)\ge0.
\]

Por construcción,

\[
D^\pi=B(E)+\Xi^\pi.
\]

La igualdad es una identidad una vez definido \(\Xi^\pi\); el resultado
sustantivo es que \(B(E)\) constituye un piso independiente del orden. Para dos ejecuciones del mismo \(E\),

\[
D^\pi-D^{\pi'}
=
\Xi^\pi-\Xi^{\pi'}.
\]

Si la red está balanceada agente por agente,

\[
\Delta(E)=0,
\]

entonces

\[
B(E)=0,
\qquad
D^\pi=\Xi^\pi.
\]

En ese caso toda necesidad distributiva dentro del horizonte considerado
es transitoria respecto de la posición terminal.

\subsection{Relevancia secuencial y testigo de cinco agentes}\label{relevancia-secuencial-y-testigo-de-cinco-agentes}

La multiplicidad de órdenes sólo es relevante para \(D\) si produce
necesidades distributivas distintas. Definamos

\[
\mathcal D(E,C)
=
\{D^\pi:\pi\in\Pi_{\mathrm{adm}}(E,C)\},
\]

\[
\mathcal X(E,C)
=
\{\Xi^\pi:\pi\in\Pi_{\mathrm{adm}}(E,C)\}.
\]

Como \(B(E)\) es común,

\[
\mathcal D(E,C)
=
B(E)+\mathcal X(E,C).
\]

La secuencia es relevante respecto de \(D\) cuando
\(|\mathcal D(E,C)|>1\). Varias cronologías pueden pertenecer a una misma
clase monetaria; la relevancia aparece cuando existe más de una.

Considérese ahora el benchmark de cinco agentes

\[
\mathcal A=(U,F,G,H_w,H_p).
\]

\(U\) produce medios de producción utilizados por \(F\) y \(G\); \(F\) y
\(G\) producen bienes consumidos por los hogares. El conjunto \(E\)
contiene doce obligaciones de salarios, consumo, reposición y
distribución del excedente. Para que el testigo sea reproducible, las
fijamos explícitamente:

\[
\begin{array}{lll}
s_U:U\to H_w,\;4, & s_F:F\to H_w,\;4, & s_G:G\to H_w,\;4,\\
c_{wF}:H_w\to F,\;6, & c_{wG}:H_w\to G,\;6, & \\[-1mm]
c_{pF}:H_p\to F,\;4, & c_{pG}:H_p\to G,\;2, & \\[-1mm]
r_F:F\to U,\;4, & r_G:G\to U,\;2, & \\[-1mm]
x_U:U\to H_p,\;2, & x_F:F\to H_p,\;2, & x_G:G\to H_p,\;2.
\end{array}
\]

Su volumen nominal es

\[
T=\sum_{e\in E}w_e=42.
\]

La red está balanceada agente por agente:

\[
\Delta(E)=0,
\qquad
B(E)=0.
\]

Para el testigo fijamos explícitamente el dominio de referencia

\[
C_{\mathrm{ref}}
=
(\varnothing,\Omega(E)),
\]

con las doce obligaciones distinguibles e indivisibles. Este dominio cumple una
función de laboratorio: no afirma que toda permutación sea contractual, jurídica o
productivamente disponible en una economía concreta. Restricciones institucionales
adicionales pueden excluir algunas comparaciones mediante \(C\). Bajo el dominio de
referencia, las dos ejecuciones que siguen pertenecen al mismo conjunto admisible.

La primera es

\[
\alpha
=
(s_U,s_F,s_G,c_{wF},c_{wG},c_{pF},c_{pG},r_F,r_G,x_U,x_F,x_G),
\]

y la segunda

\[
\beta
=
(c_{pF},c_{pG},s_U,s_F,s_G,c_{wF},c_{wG},r_F,r_G,x_U,x_F,x_G).
\]

Los mínimos de sus trayectorias son

\[
\min_k c^\alpha(k)
=
(-4,-4,-4,0,-6),
\]

\[
\min_k c^\beta(k)
=
(-4,0,-2,0,-6).
\]

Por tanto,

\[
D^\alpha=(4,4,4,0,6),
\qquad
D^\beta=(4,0,2,0,6).
\]

Como \(B(E)=0\),

\[
\Xi^\alpha=D^\alpha,
\qquad
\Xi^\beta=D^\beta.
\]

La diferencia es

\[
D^\alpha-D^\beta
=
(0,4,2,0,0).
\]

En \(\beta\), \(F\) recibe cuatro unidades y \(G\) dos antes de pagar
salarios; eso elimina la posición negativa de \(F\) y reduce la de
\(G\). Ninguna obligación cambia y el estado terminal continúa siendo
cero.

El testigo establece, por tanto, un sistema \((E,C)\) con dos ejecuciones
admisibles del mismo \(E\), igual posición terminal y distinto \(D^\pi\).
Las obligaciones completas y el cierre no determinan unívocamente la
necesidad distributiva máxima.

La diferencia entre ambas ejecuciones puede localizarse sin introducir
nuevas obligaciones. Una cadena concreta de diez transposiciones adyacentes
lleva de \(\alpha\) a \(\beta\): primero \(c_{pF}\) cruza sucesivamente
\(c_{wG},c_{wF},s_G,s_F,s_U\); después \(c_{pG}\) cruza
\(c_{wG},c_{wF},s_G,s_F,s_U\). Ocho de esos intercambios dejan \(D\)
inalterado. Sólo dos modifican un mínimo relevante: cuando \(c_{pF}\) cruza
\(s_F\), la necesidad de \(F\) cae de \(4\) a \(0\); cuando \(c_{pG}\)
cruza \(s_G\), la de \(G\) cae de \(4\) a \(2\). Bajo segregación completa,
la trayectoria del requerimiento es por tanto

\[
18\longrightarrow14\longrightarrow12,
\]

con ocho intercambios monetariamente neutrales entre esos cambios. El
resultado \(18-12=6\) puede descomponerse así en dos variaciones locales,
\(4+2\), sin atribuir efecto monetario a los intercambios que no alteran
los mínimos pertinentes.

\subsection{De la trayectoria individual a la arquitectura}\label{de-la-trayectoria-individual-a-la-arquitectura}

Hasta aquí, \(D^\pi\) conserva la profundidad máxima de cada trayectoria
individual. Para convertir esas posiciones en un requerimiento de
disponibilidad hace falta saber cuáles pueden consolidarse y si sus mínimos
coinciden. Éste es el problema de la arquitectura de acceso.


\cleardoublepage
\setcounter{subsection}{0}
\vspace*{6mm}\section*{II}\label{ii.-arquitecturas-de-compensaciuxf3n-y-requerimientos-de-liquidez}
\thispagestyle{plain}\markright{II · Arquitecturas de consolidación}
\phantomsection\addcontentsline{toc}{section}{II · Arquitecturas de consolidación y requerimientos de liquidez}
{\large\itshape Arquitecturas de consolidación y requerimientos de liquidez\par}\vspace{1.2ex}

\subsection*{Segregación, pooling y requerimiento localizado}\label{segregaciuxf3n-pooling-y-requerimiento-localizado}

\subsection{De agentes a cajas monetarias}\label{de-agentes-a-cajas-monetarias}

Dos agentes pueden atravesar posiciones de \(-2\) y \(+2\) en el mismo
prefijo. Si sus saldos permanecen segregados, la posición negativa requiere
cobertura propia; si pertenecen a una misma caja perfectamente consolidada,
ambas posiciones se consolidan en una posición neta de bloque. El requerimiento depende, por tanto, de la
arquitectura que agrupa esas trayectorias. La relevancia de sincronización, \emph{gross/net
settlement}, reutilización de pagos entrantes y mecanismos de ahorro de
liquidez está ampliamente documentada en sistemas de pagos (CPSS, 1997;
McAndrews y Rajan, 2000; Bech y Garratt, 2003; Kahn, McAndrews y
Roberds, 2003; Becher, Galbiati y Tudela, 2008; Galbiati y Soramäki,
2010; Hayakawa, 2018; McMahon et al., 2024).\footnote{Estas referencias documentan mecanismos de timing, reutilización,
\emph{gross/net settlement}, neteo y ahorro de liquidez. Hayakawa (2018), en
particular, descompone el efecto de una CCP sobre necesidades de liquidez en
efectos de enrutamiento central y neteo central; ese operador modifica la red
de obligaciones y no debe identificarse con la consolidación de posiciones
manteniendo fijo \(E\). Existe, sin embargo, un referente institucional cercano
para la semántica de acceso conjunto a saldos: TARGET2 permitió grupos de
\emph{aggregated liquidity}, en los que una orden podía liquidarse por encima
de la liquidez de una cuenta individual siempre que no excediera la liquidez
agregada de las cuentas del grupo (ECB, 2012). Las cuentas seguían siendo
jurídicamente distintas. La partición \(\mathcal P\) abstrae sólo este aspecto
de consolidación y no reproduce las reglas operativas de TARGET2.}

La pregunta deja entonces de ser cuánto necesita cada agente aisladamente y pasa a ser
qué posiciones pueden utilizarse conjuntamente para efectuar pagos. Representamos esas
fronteras de acceso monetario mediante una partición de los agentes.

Sea

\[
\mathcal P=\{P_1,\ldots,P_r\}\in\operatorname{Part}(\mathcal A)
\]

una partición de los agentes. Dentro de un bloque \(P\), las posiciones pueden consolidarse plenamente
durante la ejecución; entre bloques no existe esa capacidad. Por construcción,
``consolidarse plenamente'' significa únicamente que, para evaluar la factibilidad
monetaria de la ejecución, las posiciones positivas de los miembros del bloque pueden
cubrir sin fricción las posiciones negativas de otros miembros del mismo bloque. La
partición no establece por sí sola propiedad común, identidad jurídica, transferencia
de derechos, crédito intragrupo ni autoridad sobre los saldos. Cualquiera de esas
interpretaciones requiere reglas institucionales adicionales. \(\mathcal P\) representa,
por tanto, una frontera operativa de acceso y consolidación, no una identidad social de
los agentes que contiene.

Para cada bloque,

\[
c^\pi_P(k)
=
\sum_{a\in P}c^\pi_a(k),
\]

\[
d^\pi_P
=
\left[
-\min_{0\le k\le Q}c^\pi_P(k)
\right]_+.
\]

Podemos preguntar ahora qué saldos iniciales del medio de liquidación de referencia
tendrían que estar colocados en esas cajas para procesar toda la secuencia sin
financiación adicional. El requerimiento localizado de la arquitectura es

\[
M_0(\pi;\mathcal P)
=
\sum_{P\in\mathcal P}d^\pi_P.
\]

\(M_0\) es la suma mínima de desplazamientos iniciales no negativos
necesaria para que las trayectorias consolidadas de todas las cajas no
crucen por debajo de cero. Es, por tanto, un requerimiento de referencia
condicionado por la cronología y la arquitectura elegidas; no es por sí mismo
un agregado monetario observado, un volumen efectivo de crédito ni un coste
financiero. El origen, la disponibilidad efectiva y el precio de esos saldos
quedan fuera de \(M_0\) y se incorporan, cuando corresponde, en las capas
siguientes.

\subsection{Casos extremos y pérdida de sincronización}\label{casos-extremos-y-puxe9rdida-de-sincronizaciuxf3n}

Bajo segregación completa,

\[
\mathcal P_{\mathrm{sep}}
=
\{\{a\}:a\in\mathcal A\},
\]

cada caja contiene un solo agente y

\[
M_0(\pi;\mathcal P_{\mathrm{sep}})
=
\sum_aD^\pi_a
=
\|D^\pi\|_1.
\]

\Needspace{4\baselineskip}\textbf{Demostración.} Para cada singleton \(\{a\}\),
\(c^\pi_{\{a\}}=c^\pi_a\), por lo que \(d^\pi_{\{a\}}=D^\pi_a\). Sumando
se obtiene el resultado.

Bajo pooling completo,

\[
\mathcal P_{\mathrm{pool}}=\{\mathcal A\}.
\]

Si todas las operaciones son transferencias internas,

\[
\sum_a c^\pi_a(k)=0
\qquad
\forall k,
\]

de modo que

\[
M_0(\pi;\mathcal P_{\mathrm{pool}})=0.
\]

El requerimiento localizado interno es, por tanto, nulo cuando el conjunto
cerrado se trata como una sola caja. Este cero es un resultado de consolidación
interna: no afirma que la economía pueda liquidar obligaciones sin medios de pago,
ni que una infraestructura externa de liquidación carezca de requerimientos. Esos
medios, las obligaciones frente al exterior del bloque consolidado y las reglas de
acceso a la infraestructura quedan fuera de este caso extremo. Determinar los medios
monetarios y financieros efectivos de realización requiere información adicional.

Fuera de los extremos, \(D^\pi\) no contiene suficiente información para
calcular \(M_0\). Considérense las obligaciones unitarias

\[
e_1:A\to C,\qquad
e_2:A\to C,\qquad
e_3:C\to B,
\]

y el dominio de referencia que admite

\[
\alpha=(e_1,e_2,e_3),
\qquad
\beta=(e_1,e_3,e_2).
\]

En ambos casos,

\[
D^\alpha=D^\beta=(2,0,0).
\]

Para

\[
\mathcal P=\{\{A,B\},\{C\}\},
\]

las trayectorias consolidadas producen

\[
M_0(\alpha;\mathcal P)=2,
\qquad
M_0(\beta;\mathcal P)=1.
\]

Por tanto, en general no existe una función \(f_{\mathcal P}\) tal que

\[
M_0(\pi;\mathcal P)
=
f_{\mathcal P}(D^\pi)
\]

para toda ejecución admisible. La no-existencia es general, no una
negación para cada arquitectura: bajo segregación sí existe el caso especial
\(f(D)=\|D\|_1\).

\subsection{Orden de arquitecturas y ganancia de pooling}\label{orden-de-arquitecturas-y-ganancia-de-pooling}

Escribimos

\[
\mathcal P\preceq\mathcal P'
\]

si cada bloque de \(\mathcal P\) está contenido en algún bloque de
\(\mathcal P'\). Así, \(\mathcal P'\) permite débilmente más
consolidación.

Para bloques disjuntos \(P,Q\),

\[
d^\pi_{P\cup Q}
\le
d^\pi_P+d^\pi_Q.
\]

\Needspace{4\baselineskip}\textbf{Demostración.} Para todo \(k\),

\[
c^\pi_P(k)\ge\min_jc^\pi_P(j),
\qquad
c^\pi_Q(k)\ge\min_jc^\pi_Q(j).
\]

Por tanto,

\[
\min_k(c^\pi_P(k)+c^\pi_Q(k))
\ge
\min_kc^\pi_P(k)+\min_kc^\pi_Q(k).
\]

Como todas las trayectorias parten de cero, cambiar de signo produce la
desigualdad de requerimientos.

Toda coarsificación finita puede construirse mediante fusiones binarias.
De aquí,

\[
\mathcal P\preceq\mathcal P'
\Longrightarrow
M_0(\pi;\mathcal P')
\le
M_0(\pi;\mathcal P).
\]

En una red cerrada,

\[
0
=
M_0(\pi;\mathcal P_{\mathrm{pool}})
\le
M_0(\pi;\mathcal P)
\le
M_0(\pi;\mathcal P_{\mathrm{sep}})
=
\|D^\pi\|_1.
\]

Queda por determinar cuándo la reducción del requerimiento localizado es estricta. Definamos

\[
K^\pi_P
=
\arg\min_k c^\pi_P(k)
\]

y la holgura respecto del mínimo

\[
r^\pi_P(k)
=
c^\pi_P(k)-\min_jc^\pi_P(j)\ge0.
\]

La ganancia de fusionar \(P,Q\),

\[
\gamma^\pi(P,Q)
=
d^\pi_P+d^\pi_Q-d^\pi_{P\cup Q},
\]

mide exclusivamente la reducción de \(M_0\) producida por esa consolidación.
El término ``ganancia'' no introduce aquí una medida de bienestar, beneficio ni
ahorro de coste financiero. Además, satisface

\[
\gamma^\pi(P,Q)
=
\min_k\{r^\pi_P(k)+r^\pi_Q(k)\}.
\]

De modo que

\[
\gamma^\pi(P,Q)=0
\iff
K^\pi_P\cap K^\pi_Q\neq\varnothing,
\]

y

\[
\gamma^\pi(P,Q)>0
\iff
K^\pi_P\cap K^\pi_Q=\varnothing.
\]

La consolidación aprovecha así la holgura temporal relativa entre bloques:
la ganancia es estricta cuando sus mínimos no coinciden.

\subsection{Variaciones locales de secuencia y arquitectura}\label{variaciones-locales-secuencia-arquitectura}

\paragraph{Frontera bibliográfica.}
La operación de comparar secuencias mediante intercambios adyacentes tiene una
tradición propia en investigación de operaciones (Abdel-Wahab y Kameda, 1978;
Monma y Sidney, 1979; Monma, 1980), y el problema escalar de recursos acumulados
reaparece en Kellerer et al. (1998). En redes de obligaciones, Rotemberg (2011)
ya identifica consecuencias del orden de pago; Hayakawa et al. (2019) formulan
la liquidación secuencial como un problema de fondos de liquidación y obtienen
resultados locales para clases particulares de redes; Hayakawa (2020) estudia
cómo la topología condiciona la liquidez necesaria. McMahon et al. (2024) y
Astuti et al. (2026) llevan la optimización por reordenamiento a sistemas de
pagos de gran valor. En consecuencia, el objeto que sigue no se presenta como
una teoría nueva de que ``el orden importa'' ni como una introducción del
intercambio adyacente. La pregunta más estrecha es otra: dado el funcional de
requerimiento y una arquitectura arbitraria de consolidación, ¿cuál es la
variación exacta inducida por una transposición admisible y cómo cambia ese
margen cuando cambia la arquitectura? La revisión de los antecedentes citados
no proporciona, en esos términos, la fórmula general por bloque ni la identidad
mixta secuencia--arquitectura derivadas abajo; esta delimitación describe el
alcance de la revisión bibliográfica, no una prueba de prioridad universal.

La comparación entre cronologías puede reducirse primero a la modificación
mínima de un orden. Consideremos dos ejecuciones admisibles que sólo difieren
por una transposición adyacente,

\[
\pi=\omega ef\eta,
\qquad
\pi'=\omega fe\eta,
\]

y sea \(z\) el vector de saldos acumulados después de ejecutar \(\omega\).
Las dos trayectorias coinciden antes del par y vuelven a coincidir una vez
ejecutadas ambas obligaciones, pues

\[
z+g_e+g_f=z+g_f+g_e.
\]

Por tanto, difieren en un único estado intermedio:

\[
\pi:\quad z\longrightarrow z+g_e\longrightarrow z+g_e+g_f,
\]

\[
\pi':\quad z\longrightarrow z+g_f\longrightarrow z+g_f+g_e.
\]

Como las trayectorias sólo difieren en ese estado, la variación del
requerimiento puede calcularse sin reconstruirlas por completo. Fijemos una
arquitectura \(\mathcal P\). La proyección de una transferencia sobre sus
cajas es

\[
h_e^{\mathcal P}
:=
\left(\sum_{a\in P}g_{e,a}\right)_{P\in\mathcal P},
\qquad
h_{e,P}:=\sum_{a\in P}g_{e,a}.
\]

Si los extremos de \(e\) pertenecen al mismo bloque, entonces
\(h_{e,P}=0\) para todo \(P\): la transferencia es interna a esa
arquitectura. Para una transferencia entre bloques, las únicas componentes
no nulas son las cajas del pagador y del receptor.

Sea \(i=|\omega|+1\) la posición ocupada por la primera obligación del par.
Para cada bloque \(P\in\mathcal P\), sea \(z_P=c_P^\pi(i-1)\) su saldo
antes del par. Como el único estado distinto aparece en \(k=i\), definimos

\[
R_P
:=
\min_{\substack{0\le k\le Q\\ k\ne i}}c_P^\pi(k)
=
\min_{\substack{0\le k\le Q\\ k\ne i}}c_P^{\pi'}(k),
\qquad
\bar m_P=[-R_P]^+.
\]

Entonces

\[
d_P^\pi
=
\max\{\bar m_P,-z_P-h_{e,P}\},
\]

\[
d_P^{\pi'}
=
\max\{\bar m_P,-z_P-h_{f,P}\}.
\]

En consecuencia, si

\[
\delta_{ef}^{\pi}(\mathcal P)
:=
M_0(\pi';\mathcal P)-M_0(\pi;\mathcal P),
\]

se obtiene la identidad local

\[
\boxed{
\delta_{ef}^{\pi}(\mathcal P)
=
\sum_{P\in\mathcal P}
\left[
\max\{\bar m_P,-z_P-h_{f,P}\}
-
\max\{\bar m_P,-z_P-h_{e,P}\}
\right].
}
\]

\(\delta_{ef}^{\pi}(\mathcal P)\) compara localmente dos órdenes admisibles
manteniendo fijos \(E\) y \(\mathcal P\). Si es negativo, el segundo orden
requiere menos saldo inicial de referencia según \(M_0\). Nada se sigue todavía
sobre la capacidad de los agentes para escogerlo, los costes o conflictos de
modificar la cronología, ni el coste financiero efectivamente observado. Para
responder esas preguntas hacen falta reglas adicionales sobre formación de la
cronología, derechos de ejecución y financiación.

La fórmula separa una modificación del estado intermedio de su efecto
sobre el mínimo relevante. Definiendo

\[
x_P=-z_P-h_{e,P},
\qquad
y_P=-z_P-h_{f,P},
\]

la contribución \(\Delta_P=d_P^{\pi'}-d_P^\pi\) satisface

\[
\Delta_P>0
\iff
y_P>\max\{\bar m_P,x_P\},
\]

\[
\Delta_P<0
\iff
x_P>\max\{\bar m_P,y_P\},
\]

y

\[
\Delta_P=0
\iff
x_P=y_P
\quad\text{o}\quad
\max\{x_P,y_P\}\le\bar m_P.
\]

La neutralidad tiene así dos fuentes distintas: el intercambio puede no
modificar el estado monetario intermedio del bloque, o puede modificarlo
sin alcanzar un déficit más profundo que el mínimo ya fijado en otra
parte de la trayectoria. En particular,

\[
h_{e,P}=h_{f,P}
\Longrightarrow
\Delta_P=0,
\]

de modo que

\[
\operatorname{supp}(d^{\pi'}-d^\pi)
\subseteq
\operatorname{supp}(h_f^{\mathcal P}-h_e^{\mathcal P}).
\]

La sensibilidad local depende, por tanto, de la secuencia y de la
arquitectura de consolidación. Esta última determina qué transferencias
alteran la posición neta entre cajas y cuáles quedan internalizadas en un
mismo bloque.

Esta dependencia puede relacionarse exactamente con la ganancia de
consolidación obtenida arriba. Sea \(\mathcal P'\) la arquitectura resultante de
fusionar dos bloques \(P,Q\in\mathcal P\), y escribamos

\[
\gamma^\rho(\mathcal P\to\mathcal P')
:=
M_0(\rho;\mathcal P)-M_0(\rho;\mathcal P')\ge0.
\]

Como los bloques no fusionados aportan lo mismo en ambas arquitecturas,

\[
\gamma^\rho(\mathcal P\to\mathcal P')
=
d_P^\rho+d_Q^\rho-d_{P\cup Q}^\rho
=
\gamma^\rho(P,Q).
\]

Así, esta notación no introduce una segunda medida de ganancia: expresa
la misma variación arquitectónica en términos del funcional total.

Para \(\pi'=\tau_i\pi\), la definición implica

\[
M_0(\rho;\mathcal P')
=
M_0(\rho;\mathcal P)-\gamma^\rho
\]

para \(\rho\in\{\pi,\pi'\}\). Restando las dos igualdades se obtiene

\[
\boxed{
\delta_i^\pi(\mathcal P')-\delta_i^\pi(\mathcal P)
=
\gamma^\pi-\gamma^{\tau_i\pi}.
}
\]

La identidad relaciona dos diferencias finitas del mismo funcional; no
añade una hipótesis causal. Una consolidación deja inalterado el efecto local
de un intercambio sobre \(M_0\) si y sólo si ese intercambio deja inalterada
la reducción de \(M_0\) producida por la consolidación. Sin embargo, la
antitonicidad del nivel,

\[
\mathcal P\preceq\mathcal P'
\Longrightarrow
M_0(\rho;\mathcal P')\le M_0(\rho;\mathcal P),
\]

no implica antitonicidad de la sensibilidad \(|\delta_i^\pi|\). Un
contraejemplo balanceado requiere sólo tres agentes y cuatro obligaciones
unitarias. Sean

\[
a:A\to B,\qquad b:A\to C,\qquad
c:C\to A,\qquad d:B\to A,
\]

con

\[
\pi=(a,b,c,d),
\qquad
\pi'=(a,c,b,d).
\]

Bajo segregación completa,

\[
M_0(\pi;\mathcal P_{\mathrm{sep}})=2,
\qquad
M_0(\pi';\mathcal P_{\mathrm{sep}})=2,
\]

por lo que la diferencia local bajo segregación es cero. Si se consolidan
\(B\) y \(C\),

\[
\mathcal P'=\{\{A\},\{B,C\}\},
\]

se obtiene

\[
M_0(\pi;\mathcal P')=2,
\qquad
M_0(\pi';\mathcal P')=1,
\]

y la diferencia local pasa a \(-1\). La consolidación puede, pues,
crear sensibilidad secuencial donde antes no la había.

También puede eliminarla. Para el ciclo unitario

\[
a:A\to B,\qquad b:B\to C,\qquad c:C\to A,
\]

comparemos \(\pi=(a,b,c)\) con \(\pi'=(b,a,c)\). Bajo segregación,

\[
M_0(\pi;\mathcal P_{\mathrm{sep}})=1,
\qquad
M_0(\pi';\mathcal P_{\mathrm{sep}})=2,
\]

mientras que bajo \(\{\{A,B\},\{C\}\}\) ambas ejecuciones requieren
una unidad. Así, \(\delta\) pasa de \(+1\) a \(0\). La preservación se
produce exactamente cuando \(\gamma^\pi=\gamma^{\tau_i\pi}\), como
se sigue de la identidad anterior.

La ausencia de monotonicidad alcanza también al signo. Considérese el
sistema balanceado de cuatro agentes con obligaciones

\[
\begin{array}{llll}
e_1:A\to D,\;1, & e_2:C\to A,\;1, & e_3:D\to B,\;1, & e_4:A\to B,\;2,\\
e_5:B\to C,\;1, & e_6:C\to A,\;2, & e_7:B\to C,\;2. &
\end{array}
\]

Sean

\[
\pi=(e_1,e_2,e_3,e_4,e_5,e_6,e_7),
\qquad
\pi'=(e_1,e_2,e_3,e_4,e_5,e_7,e_6),
\]

y las arquitecturas

\[
\mathcal P=\{\{A\},\{B,D\},\{C\}\},
\qquad
\mathcal P'=\{\{A,B,D\},\{C\}\}.
\]

La segunda se obtiene de la primera mediante una sola fusión. Sin embargo,

\[
M_0(\pi;\mathcal P)=4,
\qquad
M_0(\pi';\mathcal P)=3,
\]

mientras

\[
M_0(\pi;\mathcal P')=2,
\qquad
M_0(\pi';\mathcal P')=3.
\]

El mismo intercambio pasa así de \(\delta=-1\) a \(\delta=+1\): una
fusión puede hacer que un reordenamiento reduzca \(M_0\) bajo una arquitectura
y lo aumente bajo otra. No hay contradicción con la antitonicidad de \(M_0\),
pues la consolidación reduce débilmente el requerimiento de cada ejecución
considerada por separado; lo que se invierte es su comparación. El efecto
monetario de una precedencia es, en este sentido, relativo a la organización
del acceso a los saldos. La partición, por sí sola, no permite ordenar
socialmente las arquitecturas ni inferir cambios en propiedad, poder de
decisión o incidencia de los costes.

Secuencia y arquitectura constituyen, por tanto, dos márgenes distintos de variación de

\[
M_0:\Pi_{\mathrm{adm}}\times\operatorname{Part}(\mathcal A)
\longrightarrow\mathbb R_+.
\]

\subsection{Arquitecturas nulas y descomposición de bloque}\label{arquitecturas-nulas-y-descomposiciuxf3n-de-bloque}

Para cualquier trayectoria y partición,

\[
M_0(\pi;\mathcal P)=0
\iff
c^\pi_P(k)\ge0
\quad
\forall P,k.
\]

Bajo cierre interno, la suma de las trayectorias de los bloques es cero
en cada prefijo. Por tanto,

\[
M_0(\pi;\mathcal P)=0
\iff
c^\pi_P(k)=0
\quad
\forall P,k.
\]

Cuando cada paso liquida una sola obligación bilateral positiva, sea
\(G(E)\) el grafo no dirigido de obligaciones y \(\mathcal C(E)\) su
partición en componentes conexas. Entonces

\[
M_0(\pi;\mathcal P)=0
\iff
\mathcal C(E)\preceq\mathcal P.
\]

\Needspace{4\baselineskip}\textbf{Demostración.} Si cada componente de \(G(E)\) queda contenida en
un bloque, toda transferencia es intrabloque y no altera la trayectoria
consolidada. Recíprocamente, si el requerimiento es cero, las
trayectorias de bloque permanecen en cero tras cada paso; una obligación
que cruzara bloques alteraría simultáneamente la posición de dos cajas y
produciría una contradicción.

La separación entre terminal y transitorio también puede hacerse por
bloque. Sea

\[
\Delta_P(E)=\sum_{a\in P}\Delta_a(E),
\qquad
B_P(E)=[-\Delta_P(E)]_+,
\]

\[
\Xi^\pi_P=d^\pi_P-B_P(E)\ge0.
\]

Sumando,

\[
B(E;\mathcal P)
=
\sum_{P\in\mathcal P}B_P(E),
\]

\[
\Xi(\pi;\mathcal P)
=
\sum_{P\in\mathcal P}\Xi^\pi_P,
\]

y por identidad,

\[
M_0(\pi;\mathcal P)
=
B(E;\mathcal P)+\Xi(\pi;\mathcal P).
\]

El piso \(B(E;\mathcal P)\) es antítono bajo consolidación, pero \(\Xi\)
no lo es en general. Una fusión puede eliminar piso terminal y
sustituirlo por una excursión transitoria. La monotonía pertenece, por tanto, al requerimiento total y no a cada
componente por separado.

\subsection{El benchmark y la entrada al estado heredado}\label{el-benchmark-y-la-entrada-al-estado-heredado}

Para las ejecuciones del artículo I,

\[
D^\alpha=(4,4,4,0,6),
\qquad
D^\beta=(4,0,2,0,6).
\]

Bajo segregación,

\[
M_0(\alpha;\mathcal P_{\mathrm{sep}})=18,
\qquad
M_0(\beta;\mathcal P_{\mathrm{sep}})=12.
\]

Bajo pooling completo,

\[
M_0(\alpha;\mathcal P_{\mathrm{pool}})
=
M_0(\beta;\mathcal P_{\mathrm{pool}})
=
0.
\]

Para una arquitectura \(\mathcal P\) dada, \(M_0(\pi;\mathcal P)\) está
determinado por los mínimos alcanzados por las trayectorias agregadas de
sus bloques. El orden importa cuando modifica esos mínimos; un cambio de
calendario que conserve las mismas trayectorias no los modifica.

Conviene entonces separar dos variaciones:

\[
\pi\mapsto M_0(\pi;\mathcal P)
\]

cambia la ejecución con arquitectura fija, mientras

\[
\mathcal P\mapsto M_0(\pi;\mathcal P)
\]

cambia la arquitectura con ejecución fija. La primera recoge el efecto de la secuencia bajo una arquitectura dada;
la segunda, el efecto de la consolidación sobre una ejecución dada.

\(M_0\) parte todavía de una posición inicial de referencia. Para conocer
la parte efectivamente descubierta hay que incorporar los recursos con
que cada caja llega al periodo.


\cleardoublepage
\setcounter{subsection}{0}
\vspace*{6mm}\section*{III}\label{iii.-exposiciuxf3n-monetaria-y-necesidad-de-cobertura}
\thispagestyle{plain}\markright{III · Exposición monetaria}
\phantomsection\addcontentsline{toc}{section}{III · Exposición monetaria y necesidad de cobertura}
{\large\itshape Exposición monetaria y necesidad de cobertura\par}\vspace{1.2ex}

\subsection*{Estado heredado, profundidad y duración}\label{estado-heredado-profundidad-y-duraciuxf3n}

\subsection{Estado heredado y exposición residual}\label{estado-heredado-y-exposiciuxf3n-residual}

Dos cajas que recorren la misma trayectoria pueden enfrentar exposiciones
diferentes si comienzan con recursos monetarios distintos. Para determinar
esa diferencia hay que separar el estado heredado de la parte que la
institución reconoce como disponibilidad operativa.
El problema no es sólo cuánto patrimonio o riqueza existe al inicio, sino qué parte de
esa posición puede utilizarse efectivamente para pagar ahora. El estado heredado reúne
las posiciones que llegan del periodo anterior; una regla institucional selecciona de ellas
la disponibilidad operativa reconocida.

Sea

\[
h_t\in\mathcal H
\]

el estado heredado. Para una arquitectura \(\mathcal P_t\) y una
institución dada, escribimos

\[
\mu_{I,t}:\mathcal H\to\mathbb R_+^{|\mathcal P_t|},
\qquad
m_t^0=\mu_{I,t}(h_t).
\]

El subíndice recuerda que la proyección puede cambiar con la institución o
la arquitectura; así, balances distintos pueden proyectarse sobre la misma
disponibilidad y un mismo balance puede valorarse de forma distinta bajo
instituciones diferentes. No existe una identidad \(h_t=m_t^0\):
\(m_t^0=\mu_{I,t}(h_t)\) es una proyección institucional del estado heredado.

La exposición residual de la caja \(P\) es

\[
U_{P,t}(k\mid h_t)
=
\left[
-\left(m^0_{P,t}+c^\pi_{P,t}(k)\right)
\right]_+.
\]

\(U\) mide el faltante en el prefijo \(k\) antes de movilizar cobertura
adicional. Es una exposición contrafactual previa a la cobertura. La deuda
observada y el crédito efectivamente concedido se determinan en la capa de
financiación.

La distinción entre posiciones de balance, compromisos financieros y
recursos efectivamente utilizables es consistente con enfoques de
balance y consistencia stock-flujo (Minsky, 1986, 1992; Godley y Lavoie,
2007). Passarella (2023) construye un modelo que articula teoría del circuito
monetario, insumo-producto y consistencia stock-flujo; en él, la financiación
inicial y final se relaciona con préstamos y stocks financieros que se
actualizan entre periodos. Ese antecedente sitúa la combinación de
requerimiento corriente y posición heredada; la proyección institucional y
los funcionales de exposición se definen aquí.\footnote{Minsky y Godley--Lavoie aportan el contexto de
  compromisos financieros, balances e inventarios heredados. La proyección \(\mu\), la exposición \(U\), su pico \(N\) y la
  exposición-tiempo \(\Theta\) son objetos del presente programa; su
  definición corresponde a esta construcción.}

\subsection{Profundidad y duración}\label{profundidad-y-duraciuxf3n}

Una exposición puede ser intensa y breve o menos intensa y persistente. Para no
confundir ambas dimensiones necesitamos medir por separado profundidad y duración.
El máximo de exposición es

\[
N_{P,t}
=
\max_{0\le k\le Q_t}U_{P,t}(k).
\]

Para medir persistencia temporal se distinguen tiempos físicos

\[
s_{t,0}<\cdots<s_{t,Q_t},
\qquad
\Delta s_{t,k}>0,
\]

y adoptamos la convención de que \(U_{P,t}(k)\) rige durante el
intervalo posterior al prefijo \(k\). Entonces

\[
\Theta_{P,t}
=
\sum_{k=0}^{Q_t-1}
U_{P,t}(k)\Delta s_{t,k}.
\]

El pico mide profundidad; la exposición-tiempo incorpora además su
duración.

Para el testigo fijamos un dominio \(C_{\mathrm{ref}}\) que admita las
dos ejecuciones consideradas. Con intervalos unitarios y \(m_P^0=0\),
pueden obtenerse

\[
U^\alpha_P=(0,2,2,0),
\qquad
U^\beta_P=(0,2,0,0).
\]

En ambos casos,

\[
N_P=2,
\]

pero

\[
\Theta^\alpha_P=4,
\qquad
\Theta^\beta_P=2.
\]

La duración puede cambiar aun cuando se conserve el orden relativo de las
operaciones que afectan directamente a \(P\): una operación global
intercalada puede retrasar el cobro que cierra la exposición local.

\subsection{Puente con el requerimiento localizado}\label{puente-con-el-requerimiento-localizado}

El requerimiento de bloque del artículo II es

\[
d^\pi_{P,t}
=
\left[
-\min_k c^\pi_{P,t}(k)
\right]_+.
\]

Como \(c^\pi_{P,t}(0)=0\), el mínimo es no positivo. Por tanto,

\[
N_{P,t}
=
[d^\pi_{P,t}-m^0_{P,t}]_+.
\]

\Needspace{4\baselineskip}\textbf{Demostración.} Sea \(m_P=\min_k c^\pi_{P,t}(k)\le0\). Entonces

\[
N_P
=
\max_k[-(m^0_P+c^\pi_P(k))]_+
=
[-m^0_P-m_P]_+.
\]

Como \(d^\pi_P=-m_P\), se obtiene el resultado.

En consecuencia,

\[
m^0_{P,t}\ge d^\pi_{P,t}
\iff
N_{P,t}=0
\iff
U_{P,t}(k)=0\quad\forall k.
\]

La relación separa el requerimiento desde saldo inicial de referencia
\(d^\pi_P\), la disponibilidad efectiva \(m^0_P\) y la parte máxima que
permanece descubierta \(N_P\).

La trayectoria \(U\) conserva en qué prefijos aparece cada insuficiencia y su
magnitud; \(N\) y \(\Theta\) sólo resumen aspectos de ella. El mismo \(N\) o
la misma \(\Theta\) pueden, por tanto, corresponder a perfiles diferentes.

\subsection{De la exposición a la cobertura}\label{de-la-exposiciuxf3n-a-la-cobertura}

Dado \(h_t\), la proyección institucional, la arquitectura y la trayectoria,
quedan determinados \(U_t\), \(N_t\) y \(\Theta_t\). Falta todavía determinar
cómo se cubre esa exposición: cuánto falta y durante cuánto tiempo no
establece quién proporciona la capacidad de pago ni bajo qué contrato.


\cleardoublepage
\setcounter{subsection}{0}
\vspace*{6mm}\section*{IV}\label{iv.-contratos-e-implementaciuxf3n-de-la-financiaciuxf3n}
\thispagestyle{plain}\markright{IV · Contratos e implementación}
\phantomsection\addcontentsline{toc}{section}{IV · Contratos e implementación de la financiación}
{\large\itshape Contratos e implementación de la financiación\par}\vspace{1.2ex}

\subsection*{Cobertura, instrumentos y tie-up/release}\label{cobertura-instrumentos-y-tie-uprelease}

\subsection{Cobertura y posición operacional}\label{cobertura-y-posiciuxf3n-operacional}

Una misma exposición puede sostenerse mediante coberturas distintas. El problema
económico es cómo se proporciona la capacidad de pago que falta para que la ejecución
pueda continuar. Sea \(x_{P,t}(k)\ge0\) la capacidad de pago adicional efectivamente
movilizada y \(\bar L_{P,t}(k)\) la capacidad máxima accesible. La posición
operacional es

\[
b_{P,t}(k)
=
m^0_{P,t}
+
c^\pi_{P,t}(k)
+
x_{P,t}(k).
\]

Bajo \(x\ge0\),

\[
b_{P,t}(k)\ge0
\iff
x_{P,t}(k)\ge U_{P,t}(k).
\]

\Needspace{4\baselineskip}\textbf{Demostración.} Si \(m^0_P+c^\pi_P(k)\ge0\), entonces
\(U_P(k)=0\) y toda \(x_P(k)\ge0\) mantiene no negativa la posición. Si
la suma es negativa, \(U_P(k)=-(m^0_P+c^\pi_P(k))\), de modo que
\(b_P(k)\ge0\) equivale exactamente a \(x_P(k)\ge U_P(k)\).

La literatura sobre circuito monetario, crédito y balances permite separar
la necesidad de pago de su financiación y sus contrapartidas (Campbell, 2002;
Graziani, 2003; Godley y Lavoie, 2007). Trigg (2021) recorre institucionalmente el paso desde contratos privados de
deuda y su clearing hasta un sistema bancario; Passarella (2023), por su
parte, distingue financiación inicial y financiación final, conectadas con
préstamos y depósitos dentro de un marco stock-flow consistent. Ambos antecedentes sitúan mediaciones financieras concretas entre la
necesidad corriente y las posiciones de balance. Aquí distinguimos la
cobertura, entendida como servicio de liquidez, de la estructura financiera
que la implementa.\footnote{El apoyo bibliográfico
  delimita el problema institucional de financiación, crédito y
  contrapartidas. Los operadores \(\mathcal X(\mathcal K)\),
  \(\mathcal Y(U,\mathcal K)\), \(V\), \(\Psi\) y
  \(\Lambda_{\mathcal I}\) pertenecen a la construcción formal del
  programa.}

Con límites de capacidad,

\[
U_t(k)\preceq x_t(k)\preceq\bar L_t(k).
\]

La cobertura-tiempo es

\[
\mathcal T_x(x)
=
\sum_{P\in\mathcal P_t}
\sum_{k=0}^{Q_t-1}
x_P(k)\Delta s_{t,k}.
\]

Toda cobertura operacionalmente factible satisface

\[
\max_k x_P(k)\ge N_P,
\]

y

\[
\mathcal T_x(x)\ge\Theta.
\]

Estas desigualdades fijan cotas para cualquier trayectoria que cubra la
exposición. Para pasar de ellas al crédito observado y al coste financiero
hay que especificar la cobertura, el contrato y sus precios.

\subsection{Benchmark justo-a-tiempo y régimen contractual}\label{benchmark-justo-a-tiempo-y-ruxe9gimen-contractual}

La trayectoria

\[
x^{JIT}(k)=U(k)
\]

es el elemento mínimo del dominio puramente operacional bajo el orden
componente a componente. Si existe además un límite de capacidad,
\(x^{JIT}\) pertenece al dominio sólo cuando

\[
\bar L(k)\succeq U(k)
\qquad
\forall k.
\]

Y si se introduce un régimen contractual \(\mathcal K\), JIT es
admisible únicamente cuando

\[
U\in\mathcal X(\mathcal K).
\]

JIT es, por tanto, el benchmark operacional. Su admisibilidad como
contrato efectivo depende del régimen \(\mathcal K\).

No toda trayectoria de financiación técnicamente imaginable puede contratarse o
mantenerse bajo cualquier institución. Para distinguir la necesidad operacional de las
restricciones contractuales, Sea \(\mathfrak K_t\) el universo de regímenes contractuales bien
tipados y

\[
\mathcal X(\mathcal K)
\]

el conjunto de coberturas compatibles con \(\mathcal K\). La
factibilidad efectiva es

\[
\mathcal Y(U,\mathcal K)
=
\{x\in\mathcal X(\mathcal K):x(k)\succeq U(k)\ \forall k\}.
\]

Definimos

\[
V(U,\mathcal K)
=
\inf_{x\in\mathcal Y(U,\mathcal K)}
\mathcal T_x(x),
\]

con \(V=+\infty\) cuando el conjunto es vacío. Cuando \(V<\infty\),

\[
\Psi(U,\mathcal K)
=
V(U,\mathcal K)-\Theta(U)\ge0.
\]

No se supone que el ínfimo se alcance salvo que el régimen lo garantice.

Un ejemplo de vencimiento hace visible la diferencia. Sea

\[
U=(0,2,1,0)
\]

con intervalos unitarios. El benchmark JIT tiene \(\mathcal T_x=3\).
Considérese un régimen \(\mathcal K_B\) que, cuando moviliza un bloque
de dos unidades en \(k=1\), obliga a mantenerlo íntegro hasta \(k=3\) y
permite liberarlo en ese momento. La cobertura mínima permitida es

\[
x_B=(0,2,2,0),
\]

por lo que

\[
V(U,\mathcal K_B)=4,
\qquad
\Psi(U,\mathcal K_B)=1.
\]

La exposición fija un piso operacional, pero el régimen contractual
determina qué coberturas de ese piso son admisibles.

Si

\[
\mathcal X(\mathcal K_1)\supseteq\mathcal X(\mathcal K_2),
\]

entonces

\[
\mathcal Y(U,\mathcal K_1)\supseteq\mathcal Y(U,\mathcal K_2)
\]

y

\[
V(U,\mathcal K_1)
\le
V(U,\mathcal K_2).
\]

Para valores finitos,

\[
\Psi(U,\mathcal K_1)
\le
\Psi(U,\mathcal K_2).
\]

Una capacidad constante \(\bar L_P\) es suficiente, en ausencia de otras
restricciones contractuales, exactamente cuando

\[
\bar L_P\ge N_P.
\]

\subsection{Instrumentos y contrapartidas}\label{instrumentos-y-contrapartidas}

Cubrir una exposición no determina todavía qué relación financiera presta ese servicio.
La misma capacidad de pago puede provenir de una transferencia de saldos existentes, de
crédito comercial, de un préstamo bancario o de creación de depósitos, y cada forma deja
contrapartidas distintas. Por ello conviene separar el servicio de liquidez de su
implementación financiera.

Fijada una cobertura \(x\), queda por determinar qué instrumento o
cartera la implementa. La cobertura describe el servicio de liquidez; la
implementación especifica las posiciones financieras que lo proporcionan. Sea

\[
f_t(k)
=
(f_{1,t}(k),\ldots,f_{J,t}(k))
\]

un vector de posiciones financieras y

\[
x_t(k)
=
\Lambda_{\mathcal I}(f_t(k))
\]

la regla institucional de implementación.

Si existen \(f^{(1)}\ne f^{(2)}\) en el mismo dominio institucional
tales que

\[
\Lambda_{\mathrm{impl}}(f^{(1)})
=
\Lambda_{\mathrm{impl}}(f^{(2)})
=
x,
\]

la misma cobertura admite más de una cartera implementadora.

La multiplicidad tiene contenido monetario. Si un financiador transfiere
liquidez preexistente por una cantidad \(q\), la posición monetaria
conjunta de financiador y financiado no aumenta:

\[
\Delta(m_F+m_D)=0.
\]

Si una institución origina un préstamo creando simultáneamente un
depósito aceptado como medio de pago (Graziani, 2003; Lavoie, 2014;
Wray, 2015; Mitchell, Wray y Watts, 2019), una medida monetaria que
incluya ese depósito puede registrar

\[
\Delta M(d)=q.
\]

Ambos mecanismos pueden proporcionar al financiado la misma capacidad de
pago \(q\) y el mismo pasivo principal. La cobertura observada requiere, por tanto, información adicional sobre el
mecanismo de creación o transferencia monetaria.

El crédito comercial requiere otra delimitación. Si el vendedor pospone la
liquidación de una obligación originalmente pagadera en un prefijo dado,
puede cambiar \(E\), \(C\) o la cronología efectiva. El diferimiento cuenta
como cobertura de una ejecución comercial fija sólo mientras conserve esos
tres objetos; si altera alguno, pertenece a una construcción distinta.

\subsection{Tie-up, release y estado terminal}\label{tie-up-release-y-estado-terminal}

La terminología de \emph{release} y \emph{tie-up} reúne antecedentes de
tradiciones distintas: Marx analiza liberación y vinculación de capital;
Maldonado-Filho desarrolla esas categorías en una lectura temporal; Lowe
estudia la movilización de capital circulante durante el \emph{traverse} (Marx, 1976 {[}1894{]}, cap. 6, sec.~II;
Maldonado-Filho, 1995; Lowe, 1976, cap. 17).\footnote{Las tres líneas
  cumplen funciones distintas: Marx define liberación y vinculación de
  capital; Maldonado-Filho desarrolla esas categorías en una
  reconstrucción temporal; Lowe analiza working capital dentro del
  traverse. Aquí se usan como antecedentes diferenciados, sin postular una
  genealogía Lowe \(\to\) Freeman. Los operadores \(\tau\) y \(\rho\)
  son definiciones del programa para incrementos y liberaciones de
  cobertura; la relación precisa con las categorías de Marx constituye
  una cuestión interpretativa separada.}

Para una cobertura dada, definimos

\[
\tau_P(k)
=
[x_P(k)-x_P(k-1)]_+,
\]

\[
\rho_P(k)
=
[x_P(k-1)-x_P(k)]_+.
\]

Entonces

\[
x_P(k)-x_P(k-1)
=
\tau_P(k)-\rho_P(k).
\]

La igualdad es una identidad incremental: \emph{tie-up} y \emph{release}
resumen los incrementos y liberaciones de una trayectoria de cobertura ya
dada. Bajo JIT siguen exactamente los cambios de \(U\); bajo restricciones
contractuales, una caída de la exposición puede no traducirse en liberación
inmediata.

Minimizar la necesidad de liquidez o su cobertura no basta para ordenar
institucionalmente las implementaciones. Éstas pueden diferir en robustez
frente a retrasos, impagos, concentración de dependencias o perturbaciones,
aun cuando una requiera menos liquidez.
Evaluar esas propiedades exige objetos adicionales a \(M_0\), \(N\),
\(\Theta\) y \(V\).

La cobertura describe cómo se sostuvo la ejecución presente. Para conocer
qué queda al final hay que registrar además las posiciones persistentes
producidas por su implementación. Dos implementaciones que prestan el mismo
servicio durante el periodo pueden dejar depósitos, préstamos, vencimientos
y contrapartidas diferentes. El artículo V estudia qué parte de esas
diferencias sobrevive al cierre.


\cleardoublepage
\setcounter{subsection}{0}
\vspace*{6mm}\section*{V}\label{v.-estados-financieros-frontera-y-reproducciuxf3n-interperiuxf3dica}
\thispagestyle{plain}\markright{V · Estados financieros y frontera}
\phantomsection\addcontentsline{toc}{section}{V · Estados financieros, frontera y reproducción interperiódica}
{\large\itshape Estados financieros, frontera y reproducción interperiódica\par}\vspace{1.2ex}

\subsection*{Persistencia, herencia y memoria monetaria}\label{persistencia-herencia-y-memoria-monetaria}

\subsection{Cobertura presente y estado terminal}\label{cobertura-presente-y-estado-terminal}

La cobertura explica cómo pudo continuar la ejecución; no dice todavía qué balance
monetario-financiero fue dejando esa historia. Necesitamos, por tanto, registrar depósitos,
préstamos y otras posiciones persistentes generadas durante la implementación.

Sea

\[
Z_t^M(k)\in\mathcal Z^M
\]

el estado monetario-financiero relevante en el prefijo \(k\). Su contenido depende de la institución: puede incluir caja, depósitos,
préstamos, reservas, vencimientos y posiciones frente a terceros.

Una trayectoria de cobertura no basta para describir el balance que deja la
ejecución. Dos historias financieras pueden producir el mismo \(x_t(k)\) y
acumular componentes distintos de \(Z_t^M(k)\).

En particular,

\[
x\nRightarrow Z_t^M(Q_t)
\]

en general. La cobertura identifica el servicio de liquidez prestado durante la ejecución,
pero no el estado terminal: dentro de un mismo dominio institucional, una misma
cobertura puede admitir implementaciones que dejan balances diferentes.

Esta apertura del estado a posiciones de balance y compromisos persistentes
tiene antecedentes en Minsky (1986, 1992), Godley y Lavoie (2007) y en la
articulación circuito--insumo-producto--stock-flujo de Passarella (2023),
quien construye una dinámica con dinero crediticio endógeno y hace explícita
la actualización interperiódica de stocks mediante flujos corrientes. Estos antecedentes
fundamentan el paso flujo--stock--periodo siguiente, sin definir el estado
\(Z_t^M\), la regla \(\Phi_M\) ni las preguntas de identificación
utilizadas aquí.

El contraste entre transferencia de fondos preexistentes y crédito
creador de depósito proporciona un testigo cuando ambas implementaciones
pertenecen a la misma fibra de cobertura y se mantienen fijas las demás
condiciones pertinentes.

\subsection{Regla de frontera y no-propagación}\label{regla-de-frontera-y-no-propagaciuxf3n}

El paso del estado terminal al punto de partida del periodo siguiente
requiere especificar qué parte sobrevive. Sea

\[
\Phi_M:\mathcal Z^M\times\mathcal J\to\mathcal H
\]

una regla institucional de frontera; entonces

\[
h_{t+1}
=
\Phi_M(Z_t^M(Q_t),J_t).
\]

\(J_t\) reúne información institucional no contenida necesariamente en
el vector terminal: vencimientos, reglas de liquidación, rollover y
otras condiciones.

El mismo estado terminal puede producir herencias distintas si cambia la
regla o el parámetro institucional. Por ejemplo, para el estado

\[
Z=(4,2),
\]

una regla revolvente puede producir

\[
\Phi_M^{\mathrm{roll}}(Z)=(4,2),
\]

mientras una regla que obliga a liquidar las dos unidades de deuda con
caja puede producir

\[
\Phi_M^{\mathrm{liq}}(Z)=(2,0).
\]

La diferencia procede aquí de la regla institucional: fijada \(\Phi_M\),
cada argumento tiene una imagen unívoca.

Una regla de frontera también puede borrar diferencias entre estados terminales.
Por ello,

\[
x\nRightarrow Z_t^M(Q_t)
\]

no implica automáticamente

\[
x\nRightarrow h_{t+1}.
\]

La no-identificación se propaga sólo cuando la transformación posterior conserva
las diferencias relevantes.

\subsection{T2-M: cierre comercial y memoria financiera}\label{t2-m-cierre-comercial-y-memoria-financiera}

El benchmark de cinco agentes está balanceado agente por agente. Para
toda ejecución completa de la misma red,

\[
c^\pi(T)=0.
\]

Si \(S_t^{\mathrm{com}}\) representa la dotación comercial heredada,
entonces

\[
S_T^{\mathrm{com}}
=
S_0^{\mathrm{com}}.
\]

El cierre comercial deja abierta la posición financiera.

Considérese un benchmark en el que un financiador externo cubre
exactamente el déficit incremental, no hay amortización intraperiódica
de la dotación financiada y la frontera preserva bajo rollover esa
dotación comercial no negativa.\footnote{El cálculo numérico T2-M se
  restringe al par de ejecuciones de referencia \(\alpha,\beta\) y a las
  reglas de financiación y frontera especificadas.} La necesidad
incremental es

\[
D^{\mathrm{inc}}
(S_t^{\mathrm{com}},\pi)
=
[D^\pi-S_t^{\mathrm{com}}]_+.
\]

Bajo una frontera de liquidación financiera, el saldo comercial
reconstituido permanece y el pasivo se extingue:

\[
\Phi_M^{\mathrm{liq}}
:
(S_T^{\mathrm{com}},L_T)
\mapsto
(S_T^{\mathrm{com}},0).
\]

Bajo la regla revolvente del benchmark, la dotación financiada se
conserva y la actualización puede escribirse

\[
S_{t+1}^{\mathrm{com}}
=
\max(S_t^{\mathrm{com}},D^\pi)
\]

componente a componente, bajo las hipótesis indicadas.

Con

\[
D^\alpha=(4,4,4,0,6),
\qquad
D^\beta=(4,0,2,0,6),
\]

una historia que ya dejó \(S_t^{\mathrm{com}}=D^\beta\) no requiere
financiación incremental para repetir \(\beta\), mientras el paso a
\(\alpha\) exige

\[
D^{\mathrm{inc}}(D^\beta,\alpha)
=
(0,4,2,0,0).
\]

En T2-M, la memoria monetaria depende del par formado por un estado alcanzable
y una regla institucional concreta.

\subsection{Traverse financiero y límite del cierre}\label{traverse-financiero-y-luxedmite-del-cierre}

I--V cierran así una dinámica monetario-financiera condicionada: ejecución,
requerimiento, exposición, cobertura e implementación conducen a un estado
terminal, y la regla de frontera determina su herencia.

Ese cierre todavía no determina la transición material de la economía. Stocks
físicos, producción en proceso, capacidad productiva y composición sectorial
pueden cambiar entre periodos sin estar contenidos en \(h_{t+1}\). El problema de articular reproducción material,
circulación, capital circulante y cambio estructural remite a líneas
distintas en Marx, Lowe y Pasinetti (Marx, 1976 {[}1885{]}; Marx, 1976
{[}1894{]}; Lowe, 1976, caps. 12 y 17; Pasinetti, 1981).\footnote{Estas
  fuentes motivan dimensiones distintas del problema: reproducción y
  circulación en Marx, traverse y capital circulante en Lowe, cambio
  estructural interperiódico en Pasinetti. Ninguna proporciona la regla
  \(\Phi_M\) ni la transición intraperiódica del programa.}

La cadena I--V deja además un problema común. En varios pasos, dos objetos
indistinguibles para una descripción producen valores distintos del observable
siguiente. Antes de ampliar el estado hacia una dinámica conjunta conviene
preguntar qué información debe conservar una descripción para determinar cada
observable ya obtenido.


\cleardoublepage
\setcounter{subsection}{0}
\vspace*{6mm}\section*{VI}\label{vi.-resoluciuxf3n-informacional-de-sistemas-secuenciales}
\thispagestyle{plain}\markright{VI · Resolución informacional}
\phantomsection\addcontentsline{toc}{section}{VI · Resolución informacional de sistemas secuenciales}
{\large\itshape Resolución informacional de sistemas secuenciales\par}\vspace{1.2ex}

\subsection*{Fibras, cocientes y geometría de la ejecución}\label{fibras-cocientes-y-geometruxeda-de-la-ejecuciuxf3n}

\subsection{Del caso económico a la fibra}\label{del-caso-econuxf3mico-a-la-fibra}

El problema económico de este artículo puede formularse antes de introducir su
instrumental matemático: si observamos o conservamos cierta información sobre una
ejecución, ¿basta para conocer el resultado que nos interesa, o esa descripción ha
identificado casos que económicamente deben distinguirse? Las fibras y la factorización
se introducen para responder exactamente esa pregunta de suficiencia informacional.

Los artículos anteriores han encontrado repetidamente el mismo problema:
una descripción suficiente para una pregunta puede dejar de serlo para la
siguiente.

\begin{itemize}
\tightlist
\item
  dos trayectorias pueden compartir \(D\) y producir diferente \(M_0\)
  bajo una arquitectura;
\item
  una misma exposición puede admitir varias coberturas;
\item
  una misma cobertura puede admitir varias implementaciones;
\item
  una misma cobertura puede conducir a estados terminales distintos.
\end{itemize}

La cuestión que dejan planteada los resultados de I--V puede formularse ahora de manera general: dado un observable objetivo, ¿qué distinciones entre ejecuciones o mediadores deben conservarse para reconstruirlo?

Sea

\[
F:\Pi_{\mathrm{adm}}\to Y.
\]

Definimos

\[
\pi\sim_F\pi'
\iff
F(\pi)=F(\pi').
\]

Las fibras de \(F\) agrupan ejecuciones indistinguibles para esa
pregunta. El cociente

\[
\Pi_{\mathrm{adm}}/\sim_F
\]

conserva exactamente las distinciones visibles para \(F\).

Los antecedentes movilizados en I--V operan también a resoluciones
distintas: Leontief--Bródy complementa la estructura productiva con stocks
y flujos monetarios; Trigg explicita contratos, compensación y mediaciones
bancarias; Rotemberg y Hayakawa sitúan la liquidación secuencial dentro de
redes de obligaciones, mientras Hayakawa et al. y McMahon et al. convierten el
orden de pagos en objeto explícito de optimización; Passarella articula la
secuencia monetaria intraperiódica con stocks y dinámica interperiódica.
Aquí estos antecedentes se ordenan según la resolución informacional que
exige cada problema. El criterio de fibras formaliza esa comparación dentro
del presente programa.

\subsection{Teorema de factorización y criterio de identificación}\label{teorema-de-factorizaciuxf3n-y-criterio-de-identificaciuxf3n}

La pregunta anterior tiene una respuesta exacta. Una descripción determina
otra si ésta es constante sobre cada clase que la primera no distingue. En
lenguaje matemático, una función constante sobre clases de equivalencia
factoriza a través del cociente correspondiente (Leinster, 2014).\footnote{La propiedad universal del
  cociente en conjuntos proporciona el soporte matemático de la
  factorización. La lectura económica en términos de suficiencia de una
  representación para un observable es propia del programa y el resultado
  se formula aquí en términos deterministas.}

Sea

\[
F:\Pi_{\mathrm{adm}}\to Y,
\qquad
G:\Pi_{\mathrm{adm}}\to Z.
\]

Son equivalentes:

\begin{enumerate}
\def\labelenumi{\arabic{enumi}.}
\tightlist
\item
  \(\sim_F\subseteq\sim_G\);
\item
  \(G\) es constante sobre cada fibra de \(F\);
\item
  existe una única función
\end{enumerate}

\[
\bar g:\operatorname{im}F\to\operatorname{im}G
\]

tal que

\[
G=\bar g\circ F.
\]

\Needspace{4\baselineskip}\textbf{Demostración.} La inclusión de relaciones de equivalencia
significa exactamente que \(F(\pi)=F(\pi')\) implica \(G(\pi)=G(\pi')\),
que es la constancia de \(G\) en las fibras de \(F\). Bajo esa
condición, definimos

\[
\bar g(F(\pi))=G(\pi).
\]

La constancia en fibras hace que la definición sea independiente del
representante. Por construcción, \(G=\bar g\circ F\). La unicidad sobre
\(\operatorname{im}F\) es forzosa. La recíproca se obtiene aplicando
\(\bar g\) a dos valores iguales de \(F\).

Escribiremos

\[
F\succeq_{\mathrm{info}}G
\]

cuando \(G\) factoriza a través de \(F\). En términos económicos, \(F\)
conserva entonces todas las distinciones que \(G\) necesita: nunca reúne
dos casos que \(G\) deba separar. El orden resultante es determinista y
se define por refinamiento de fibras; la comparación conceptual con los
órdenes de \emph{informativeness} de Blackwell (1951) sitúa esta relación
de suficiencia.

El cociente canónico

\[
q_F:
\Pi_{\mathrm{adm}}
\to
\Pi_{\mathrm{adm}}/\sim_F
\]

es, salvo renombramiento de clases, el descriptor exacto menos
informativo que todavía determina \(F\).

El mismo criterio se aplica a mediadores institucionales. Si
\(\Gamma(X)\) es la fibra de mediadores compatibles con un objeto \(X\)
y

\[
F_X:\Gamma(X)\to Y
\]

es el resultado inducido manteniendo fijas las demás condiciones
pertinentes, entonces \(X\) identifica unívocamente \(Y\) a través de
esa mediación si y sólo si \(F_X\) es constante sobre \(\Gamma(X)\). Un
par de mediadores compatibles con resultados distintos basta para
demostrar no-identificación. Seleccionar el mediador e identificar el mecanismo que lo genera exige una
especificación institucional adicional.

\subsection{Testigos, recuperación de suficiencia y jerarquías}\label{testigos-recuperaciuxf3n-de-suficiencia-y-jerarquuxedas}

El criterio devuelve inmediatamente los resultados anteriores. Si

\[
F(\pi)=F(\pi'),
\qquad
G(\pi)\ne G(\pi'),
\]

entonces \(F\not\succeq_{\mathrm{info}}G\). Los testigos de I--V son
precisamente pares de este tipo.

La suficiencia se evalúa en cada transformación y dominio. Una
transformación posterior puede borrar diferencias que \(F\) no distinguía,
y la suficiencia puede recuperarse al restringir el dominio o fijar una
institución. Bajo segregación,

\[
M_0(\pi;\mathcal P_{\mathrm{sep}})
=
\|D^\pi\|_1,
\]

de modo que \(D^\pi\) es suficiente para ese requerimiento. Si un
régimen contractual fija una regla de selección

\[
x=\sigma_{\mathcal K}(U),
\]

entonces \(U\) basta para \(x\) en ese dominio. Fijados el dominio institucional y el mapa que lleva cada implementación a
su cobertura, la cobertura identifica unívocamente la implementación
exactamente cuando ese mapa es inyectivo.

Bajo parámetros institucionales fijos pueden establecerse jerarquías
como

\[
c^\pi\succeq_{\mathrm{info}}D^\pi.
\]

Cambiar \(\mathcal P\), \(h\), \(\mathcal K\) o \(\Phi_M\) cambia
la pregunta y puede cambiar las fibras. La clasificación informacional de
los observables queda así indexada por el dominio y las condiciones
institucionales fijadas.

\subsection{\texorpdfstring{Estados parciales, órdenes completos y laboratorio \(2\times2\)}{Estados parciales, órdenes completos y laboratorio 2\textbackslash times2}}\label{estados-parciales-uxf3rdenes-completos-y-laboratorio-2times2}

Cuando \(n\) operaciones distinguibles pueden estar ejecutadas o
pendientes, los estados parciales forman

\[
\mathcal U_n=\{0,1\}^n.
\]

Dos estados son adyacentes cuando difieren en una coordenada, dando el
hipercubo \(Q_n\). Una permutación completa corresponde a un camino
monotónico maximal de \(0^n\) a \(1^n\).

Las ejecuciones completas forman otro objeto:

\[
G_n
=
\operatorname{Cay}
(S_n,\{s_1,\ldots,s_{n-1}\}),
\]

donde \(s_i\) intercambia posiciones consecutivas. Los vértices de
\(Q_n\) son estados parciales; los vértices de \(G_n\) son órdenes
completos.

En el laboratorio de cuatro agentes \((C,V,K,W)\), con operaciones

\[
a:K\to C(4),\quad
b:C\to W(3),\quad
c:K\to V(2),\quad
d:V\to W(5),
\]

se obtiene, para toda \(\pi\in S_4\),

\[
D^\pi
=
\left(
3\mathbf1_{\{b\prec_\pi a\}},
\,
3+2\mathbf1_{\{d\prec_\pi c\}},
\,
6,
\,
0
\right).
\]

Por tanto,

\[
q_D(\pi)
=
\left(
\mathbf1_{\{b\prec_\pi a\}},
\mathbf1_{\{d\prec_\pi c\}}
\right)
\]

clasifica exactamente las cuatro clases de \(D\). Cada clase contiene
seis ejecuciones.

Para este observable y este dominio bastan, por tanto, dos comparaciones
ordinales. Cambiar el observable o el dominio puede cambiar la reducción
suficiente.

\subsection{Sensibilidad local, geometría y factibilidad}\label{sensibilidad-local-geometruxeda-y-factibilidad}

Sea \(H\) un grafo de comparaciones entre operaciones. Cada orden total
orienta sus aristas; la imagen del mapa de orientación es el conjunto
\(AO(H)\) de orientaciones acíclicas. Para bosques, toda orientación es
acíclica y la estructura de orientaciones es hipercúbica. Cuando existen
ciclos, aparecen restricciones de realizabilidad. La teoría de
orientaciones acíclicas y la geometría de permutoedros y zonotopos
proporcionan el lenguaje matemático de esta representación (Stanley,
1973; Postnikov, 2009; Ziegler, 1995).\footnote{Stanley, Postnikov y
  Ziegler suministran resultados y lenguaje matemático sobre
  orientaciones acíclicas, permutoedros y zonotopos. No se les atribuye
  el significado económico de los observables ni la selección de un
  dominio admisible.} En VI, este lenguaje representa el observable sobre un dominio admisible ya
fijado.

Para un dominio admisible \(A=\Pi_{\mathrm{adm}}\), definimos el grafo
de sensibilidad \(H_F^A\) por testigos locales: \(\{e,f\}\) es una
arista cuando existe un contexto admisible en el que intercambiar
consecutivamente \(e\) y \(f\) cambia \(F\).

En el dominio completo \(S_n\), y también en dominios definidos sólo por
precedencias bajo las condiciones correspondientes, la orientación de
\(H_F^A\) puede ser suficiente para \(F\). La razón es que dos
extensiones lineales de un mismo poset pueden conectarse mediante
intercambios adyacentes de elementos incomparables; si esos intercambios
son \(F\)-invariantes, \(F\) permanece constante a lo largo del camino.

Con restricciones adicionales de factibilidad, la conclusión debe
formularse por fibra. Sea

\[
G_{\mathrm{adm}}
=
G_n[A]
\]

y

\[
A_\alpha
=
\{\pi\in A:q_{H_F^A}(\pi)=\alpha\}.
\]

Si \(F\) es invariante en cada arista de \(G_{\mathrm{adm}}[A_\alpha]\),
entonces \(F\) es constante en cada componente conexa de esa fibra. Si
la fibra inducida es conexa, \(F\) es constante en toda la fibra. La
conectividad es suficiente, no necesaria: una fibra desconectada puede
ser globalmente constante por razones adicionales.

El contraejemplo mínimo

\[
A=\{abc,cba\}
\]

muestra el problema. Las dos ejecuciones no están conectadas por
intercambios adyacentes admisibles, de modo que la sensibilidad local
puede ser vacía aun cuando \(F\) tome valores distintos. Cuando la fibra
se fragmenta, puede enriquecerse el descriptor:

\[
\widehat q_F(\pi)
=
\bigl(
q_{H_F^A}(\pi),
\operatorname{comp}(\pi)
\bigr).
\]

La información faltante puede provenir de la factibilidad, sin implicar una
interacción económica de orden superior.

La definición anterior registra si un intercambio puede cambiar un
observable. Para el observable \(M_0(\cdot;\mathcal P)\), el artículo II
permite además cuantificar esa arista: a cada transposición admisible
\(\pi\to\tau_i\pi\) puede asociarse

\[
\delta_i^\pi(\mathcal P)
=
M_0(\tau_i\pi;\mathcal P)-M_0(\pi;\mathcal P).
\]

El grafo de sensibilidad distingue entonces aristas activas y neutrales,
mientras la variación local conserva signo y magnitud. Como
\(\delta_i^\pi\) depende de \(\mathcal P\), una misma comparación puede
ser visible bajo una arquitectura y neutral bajo otra. La geometría de
las ejecuciones y el orden de arquitecturas no son, por tanto, dos
reducciones independientes del problema: son dos direcciones de
variación del mismo funcional.

\subsection{Uso operativo y suficiencia prospectiva}\label{uso-operativo-y-suficiencia-prospectiva}

Una vez obtenido un descriptor suficiente, las clases relevantes pueden
cubrirse o enumerarse sin recorrer cada ejecución por separado. Las construcciones combinatorias de cobertura cumplen esa función sobre un
dominio ya definido.\footnote{Deben distinguirse dos problemas combinatorios
  que a veces se confunden. Una superpermutación contiene cada
  permutación como subcadena contigua (Johnston, 2013; Houston, 2014),
  mientras una \emph{shortest common supersequence} exige aparición como
  subsecuencia y admite borrado de símbolos intermedios (Räihä y
  Ukkonen, 1981). Ningún resultado económico de VI depende de que una
  construcción de cobertura sea mínima; su función, cuando se usa, es
  sólo recorrer o certificar un dominio previamente definido.} La
dirección lógica es

\[
(E,C)\mapsto\Pi_{\mathrm{adm}}(E,C)
\]

y sólo después se pregunta si una construcción recorre adecuadamente ese
dominio.

El mismo criterio puede aplicarse prospectivamente a la herencia: fijados
el dominio, la institución de financiación y la regla de frontera, puede
preguntarse qué observable presente basta para determinar \(h_{t+1}\) o un
estado futuro más rico.

VI responde así al problema recurrente de I--V: una representación es
suficiente para un observable exactamente cuando conserva todas las
distinciones que éste requiere. Con esta cuestión resuelta, queda abierta
otra de naturaleza económica: qué diferencias producidas dentro del
período sobreviven a la frontera y modifican el estado reproductivo
siguiente. Ése es el problema de VII.


\cleardoublepage
\setcounter{subsection}{0}
\vspace*{6mm}\section*{VII}\label{vii.-reproducciuxf3n-traverse-y-secuenciaciuxf3n-monetaria}
\thispagestyle{plain}\markright{VII · Reproducción, traverse y secuenciación}
\phantomsection\addcontentsline{toc}{section}{VII · Reproducción, traverse y secuenciación monetaria}
{\large\itshape Reproducción, traverse y secuenciación monetaria\par}\vspace{1.2ex}

\subsection*{Planteamiento de una dinámica sensible a la ejecución intraperiódica}\label{planteamiento-de-una-dinuxe1mica-de-dos-tiempos}

\subsection{Dónde termina la teoría cerrada}\label{duxf3nde-termina-la-teoruxeda-cerrada}

La teoría cerrada en V alcanza un estado heredable \(h_{t+1}\), formado
por depósitos, pasivos, liquidez y otras posiciones financieras. Una
dinámica reproductiva conjunta requiere además stocks físicos, producción
en proceso, capacidad productiva y composición sectorial. El criterio de VI permite evaluar qué descripciones conservan información
suficiente una vez definido el objeto futuro. Antes hace falta construir esas
dimensiones y relacionarlas con la regla monetario-financiera \(\Phi_M\).

Passarella (2023) construye una representación stock-flow consistent que
articula la secuencia del circuito dentro del periodo con la dinámica entre
periodos y extiende esa articulación a una economía multisectorial mediante
una estructura insumo-producto. El problema que queda abierto aquí es más
estrecho: determinar cuándo diferencias entre ejecuciones admisibles dentro
de un mismo periodo sobreviven a la frontera como diferencias del estado
reproductivo siguiente.

Para formular esa cuestión conviene distinguir dos escalas temporales:

\[
t=0,1,2,\ldots
\]

para el tiempo reproductivo/interperiódico, y

\[
k=0,\ldots,Q_t
\]

para la ejecución intraperiódica. La distinción mantiene separadas la
secuencia interna de pagos y la transición reproductiva \(t\to t+1\): el
tiempo de prefijos describe la primera, mientras la segunda exige una
construcción interperiódica adicional.

\subsection{Estado candidato y operadores abiertos}\label{estado-candidato-y-operadores-abiertos}

A diferencia de los objetos cerrados de I--V, el estado que sigue es deliberadamente
un objeto candidato: nombrarlo no sustituye la tarea económica de especificar qué
dimensión material debe conservar y cómo se relaciona con la ejecución monetaria.
Para articular ambas escalas basta, por ahora, con introducir un estado

\[
B_t(k).
\]

Una descomposición posible es

\[
B_t(k)
=
\bigl(
B_t^R(k),
B_t^M(k)
\bigr),
\]

donde \(B_t^R\) recogería la dimensión real/material y \(B_t^M\) la
monetario-financiera. Esta descomposición fija sólo el objeto que habrá que construir: quedan por
determinar el tipo de cada componente, sus invariantes y la resolución
mínima que debe conservar.

Una dinámica conjunta requeriría, al menos, una actualización
intraperiódica

\[
B_t(k)
\xrightarrow{\mathcal T_e}
B_t(k+1)
\]

y una frontera

\[
B_t(Q_t)
\xrightarrow{\Phi}
B_{t+1}(0).
\]

\(\mathcal T_e\) y \(\Phi\) nombran los operadores que una dinámica
conjunta deberá definir. Su construcción debe especificar
dominio, codominio, invariantes y condiciones de buena definición. Para recuperar la regla monetario-financiera \(\Phi_M\) de V como componente
o proyección, primero debe establecerse una relación entre ambos espacios de
estado; por ello, en la construcción abierta,

\[
\Phi\ne\Phi_M
\]

en general.

\subsection{Antecedentes y diferenciación del problema}\label{antecedentes-y-diferenciacion-del-problema}

La pregunta abierta puede situarse con mayor precisión si se distinguen
problemas que la literatura sobre reproducción fue separando progresivamente.
Los esquemas de reproducción permiten establecer relaciones de reposición y
continuidad; no determinan, por ese solo hecho, una dinámica particular.
Tugan-Baranovsky utilizó esos esquemas para examinar proporcionalidad y crisis;
Luxemburg desplazó la discusión hacia las condiciones de realización; Bauer
construyó trayectorias multiperiódicas de reproducción ampliada y Grossmann
preguntó por su sostenibilidad bajo reglas determinadas de acumulación. Feldman
convirtió posteriormente la reproducción sectorial en un problema explícito de
crecimiento y asignación de la inversión. En estos desarrollos se separan así
preguntas distintas: qué configuración es reproductivamente compatible, cómo
puede realizarse, qué reglas generan una trayectoria y si esa trayectoria puede
sostenerse (Luxemburg, 1913; Grossmann, 1929; Gehrke, 2018).

Otra línea convirtió esas relaciones en sistemas multisectoriales susceptibles
de determinación matemática. Leontief hizo explícita la interdependencia
productiva mediante coeficientes técnicos; von Neumann incorporó la expansión y
su sistema dual de precios; Seton formuló matricialmente relaciones entre
valores y precios; Morishima y Seton hicieron de la agregación en departamentos
una propiedad que debía establecerse, en lugar de tomarla como dada. Okishio
incorporó además el cambio técnico bajo criterios de selección capitalista.
Morishima pudo así reconstruir los esquemas marxianos dentro de una teoría
multisectorial de valor y crecimiento y preguntar por acumulación y estabilidad
(Morishima y Seton, 1961; Morishima, 1973). La reproducción deja aquí de
proporcionar únicamente condiciones de compatibilidad: una vez añadidas reglas
de comportamiento y transición, puede generar una trayectoria determinada.

Kalecki llega al mismo objeto desde otra dirección. Su teoría de la inversión,
los beneficios y la demanda efectiva le permite reinterpretar las ecuaciones de
reproducción como relaciones entre departamentos verticalmente integrados y
componentes de la demanda (Kalecki, 1968). El desplazamiento es relevante aquí
porque separa nuevamente dos preguntas: una estructura puede satisfacer
relaciones reproductivas sin que éstas determinen por sí mismas el gasto
mediante el cual se realiza el producto. La realización requiere un cierre
adicional.

El problema monetario introduce todavía otra mediación. En Marx, la
reproducción del capital social ya contiene circulación monetaria, rotación,
crédito y funciones del dinero; la cuestión no consiste, por ello, en añadir
dinero a una construcción originalmente no monetaria. De Brunhoff reconstruye
la articulación entre dinero, medio de pago, crédito y capital; en particular,
la separación entre entrega y pago permite que una transacción constituya una
obligación cuyo vencimiento y liquidación ocurran posteriormente (de Brunhoff,
1976 {[}1967{]}). Foley recoge directamente esa reconstrucción y plantea su
formalización. Sus trabajos de 1982 realizan dos operaciones diferentes: uno
construye una relación operacional entre magnitudes monetarias y trabajo social;
el otro formula un modelo temporal del circuito en el que producción, venta,
realización y reinversión pueden quedar separadas por duraciones y rezagos
(Foley, 1982a, 1982b). Bajo crecimiento positivo y rezagos positivos de
recompromiso del valor realizado, el segundo modelo requiere expansión del
crédito para permitir la venta de mercancías. El punto no es simplemente que el
dinero importe, sino que stocks, flujos y tiempos de circulación entran en una
representación cuantitativa del movimiento del capital.

Estas líneas permiten ordenar una secuencia de problemas sin convertirla en una
genealogía causal única:

\[
\text{reproducción}
\longrightarrow
\text{realización}
\longrightarrow
\text{financiación}
\longrightarrow
\text{liquidación}
\longrightarrow
\text{secuenciación}.
\]

Las flechas indican requerimientos adicionales de información, no una cadena de
influencias ni una causalidad automática. Una condición reproductiva no
determina su mecanismo de realización; la realización no determina cómo se
financian los desfases entre pagos y cobros; disponer de financiación no
determina la arquitectura mediante la cual se liquidan las obligaciones; y
fijar obligaciones y arquitectura tampoco determina necesariamente el orden
concreto de ejecución.

El objeto construido en I--V ocupa precisamente los últimos pasos de esta
resolución. Una estructura reproductiva \(R\) no contiene por sí sola una red
bilateral de obligaciones. La mediación institucional \(I_\chi\) permite
constituir \(E\); las restricciones pertinentes delimitan después
\(\Pi_{\mathrm{adm}}\). Mantener \(E\) fijo y variar
\(\pi\in\Pi_{\mathrm{adm}}\) permite aislar una pregunta distinta: qué
diferencias aparecen exclusivamente porque las mismas obligaciones se ejecutan
en órdenes diferentes.

Los resultados de I--V no requieren atribuir a esa secuencia un mecanismo
causal adicional. Para una ejecución dada, la trayectoria acumulada determina
una necesidad distributiva; una arquitectura de consolidación convierte esa
información en un requerimiento localizado; los recursos heredados permiten
distinguir el requerimiento de la exposición efectiva; y el régimen contractual
y la implementación financiera determinan qué coberturas son admisibles y qué
posiciones persisten al cierre. En cada paso se conserva el objeto anterior y se
incorpora la información necesaria para responder una pregunta que permanecía
indeterminada.

La secuenciación ocupa, por tanto, un lugar más estrecho que una teoría general
del dinero o del crédito. Que una obligación pueda financiarse no determina
cuándo será liquidada; que todas las obligaciones terminen pagándose tampoco
determina las posiciones transitorias recorridas durante su ejecución. Dos
órdenes pueden conservar la misma red \(E\) y el mismo resultado terminal y
producir, sin embargo, trayectorias acumuladas y requerimientos de disponibilidad
distintos. El resultado establecido hasta aquí pertenece a esa construcción:
dadas las obligaciones, la arquitectura y las demás condiciones fijadas,
modificar la ejecución puede modificar los observables que dependen de sus
prefijos.

Esta delimitación separa también el problema presente de la literatura de
sistemas de pagos. El efecto de la reutilización de entradas, la compensación,
el neteo o el orden de liquidación sobre las necesidades de liquidez constituye
un problema conocido una vez dada una red de pagos. Aquí la pregunta adicional
aparece aguas arriba y aguas abajo de ese resultado. Aguas arriba, interesa cómo
una estructura reproductiva llega a constituirse como una red bilateral de
obligaciones. Aguas abajo, interesa determinar si las diferencias producidas
durante su ejecución desaparecen al cerrar el período o sobreviven como
diferencias relevantes para la reproducción posterior.

Lowe y Pasinetti delimitan dos antecedentes adicionales para este último paso.
Lowe formula el \emph{traverse} y el problema del capital circulante durante una
transición; Pasinetti desarrolla una teoría de cambio estructural en la que los
estados heredados y la transformación interperiódica son centrales (Lowe, 1976,
caps. 12 y 17; Pasinetti, 1981). Las reconstrucciones temporales de Marx ayudan,
por otra vía, a preservar la determinación secuencial frente a lecturas
simultáneas (Freeman, 1995; Tombazos, 2014). Estas líneas aportan antecedentes
convergentes para problemas distintos; no identifican todavía la transición
\((B_t(k),\mathcal T_e,\Phi)\), que debe construirse y probarse en el nuevo
espacio de estado.

\subsection{De la secuenciación al traverse}\label{de-la-secuenciacion-al-traverse}

La cuestión fuerte puede formularse ahora con precisión: \textbf{¿bajo qué
instituciones una diferencia de ejecución dentro del período \(t\) sobrevive a
la frontera como una diferencia reproductiva en \(t+1\)?} Los resultados
anteriores permiten determinar cómo diferencias de secuencia pueden producir
diferencias de necesidad, exposición, financiación y estado
monetario-financiero terminal. No permiten concluir todavía que dos ejecuciones
distintas produzcan dos trayectorias reproductivas distintas. Para obtener ese
resultado hace falta especificar qué componentes del estado terminal atraviesan
la frontera y cómo intervienen en la transición real-material.

La distinción entre las escalas \(k\) y \(t\) adquiere aquí su contenido
económico. El índice \(k\) ordena eventos dentro de un período y permite
reconstruir la trayectoria de ejecución; el índice \(t\) ordena estados
reproductivos sucesivos. Que una diferencia exista para algún \(k<Q_t\) no
implica que sobreviva cuando comienza \(t+1\). Incluso una diferencia presente
al cierre monetario puede resultar irrelevante para la reproducción posterior
si la regla de frontera la elimina, la compensa o la proyecta sobre estados
reales equivalentes.

Por ello, el estado candidato

\[
B_t(k)=\bigl(B_t^R(k),B_t^M(k)\bigr)
\]

no constituye todavía una teoría cerrada de la dinámica conjunta. Su función es
localizar el objeto que falta. \(B_t^M(k)\) debe conservar las posiciones
monetario-financieras pertinentes; \(B_t^R(k)\) deberá representar las
dimensiones materiales que la transición reproductiva necesite distinguir. La
forma concreta de estas últimas ---stocks físicos, producción en proceso,
capacidad, inventarios o composición sectorial--- no puede decidirse antes de
especificar qué transición se pretende determinar.

La frontera conjunta

\[
B_t(Q_t)\xrightarrow{\;\Phi\;}B_{t+1}(0)
\]

es entonces el lugar en que puede aparecer o desaparecer la dependencia de
trayectoria. Si dos ejecuciones \(\pi_t\) y \(\pi_t'\) parten de la misma
estructura reproductiva y satisfacen las mismas restricciones, todavía debe
preguntarse si

\[
\Phi\!\left(B_t^{\pi_t}(Q_t)\right)
=
\Phi\!\left(B_t^{\pi_t'}(Q_t)\right)
\]

o si, bajo determinadas instituciones,

\[
\Phi\!\left(B_t^{\pi_t}(Q_t)\right)
\ne
\Phi\!\left(B_t^{\pi_t'}(Q_t)\right).
\]

Sólo el segundo caso permite afirmar que una diferencia producida por la
ejecución atravesó la frontera del período. Aun entonces debe distinguirse qué
componente explica la diferencia: una posición financiera heredada, una
restricción de financiación, un inventario no realizado, producción
interrumpida, capacidad modificada u otra variable incorporada explícitamente a
\(B^R\).

T2-M exhibe una forma de memoria monetaria bajo las reglas de \emph{rollover}
especificadas. Ese resultado no demuestra que la memoria alcance al estado
real-material. Una frontera conjunta puede borrar diferencias intraperiódicas y
hacer converger ejecuciones distintas al mismo estado futuro. La hipótesis que
VII deja abierta es más precisa: existen instituciones y reglas de transición
bajo las cuales algunas diferencias generadas por la ejecución intraperiódica
permanecen informacional y económicamente relevantes después de atravesar la
frontera.

La cadena

\[
R_t
\to E_t
\to \Pi_{\mathrm{adm}}(E_t,C_t)
\to \pi_t
\to B_t(Q_t)
\xrightarrow{\Phi} B_{t+1}(0)
\]

organiza las condiciones que una construcción dinámica deberá satisfacer.
\(E_t\) requiere una mediación institucional; la admisibilidad debe fijarse
antes de los resultados monetarios; \(B_t\), \(\mathcal T_e\) y \(\Phi\) deben
construirse explícitamente. La construcción conjunta deberá recuperar como
proyecciones los resultados cerrados de I--V cuando sus objetos de entrada
permanezcan fijos.

El problema del \emph{traverse} puede plantearse así como una cuestión de
persistencia. Una modificación de la ejecución altera primero una trayectoria
dentro del período. La arquitectura institucional determina si esa modificación
altera el requerimiento o la exposición; la financiación determina qué estado
financiero resulta; la regla de frontera determina qué componentes son
heredados; finalmente, la transición real-material determina si la diferencia
heredada modifica las condiciones de reproducción. No existe una inferencia
directa desde el primer término hasta el último.

\subsection{Pregunta informacional del estado futuro}\label{pregunta-informacional-del-estado-futuro}

Sólo una vez definido \(B_{t+1}(0)\) puede aplicarse el criterio de VI y
preguntar si un observable presente \(F\) basta para determinarlo:

\[
F\succeq_{\mathrm{info}}B_{t+1}(0)\ ?
\]

Entonces podrá preguntarse si basta \(D^\pi\), si hace falta \(U\), si debe
conservarse la trayectoria completa \(c^\pi\), la implementación financiera
o un descriptor más rico. Antes de fijar el estado futuro no hay todavía
un resultado de suficiencia que demostrar.

\subsection{Qué contaría como cierre}\label{quuxe9-contaruxeda-como-cierre}

Cerrar la dinámica exige, como mínimo:

\begin{enumerate}
\def\labelenumi{\arabic{enumi}.}
\tightlist
\item
  definir el tipo y contenido mínimo de \(B_t(k)\);
\item
  construir una transición intraperiódica \(\mathcal T_e\) bien
  definida;
\item
  construir la frontera conjunta \(\Phi\);
\item
  demostrar compatibilidad con las capas I--V;
\item
  obtener al menos un resultado de dependencia o independencia de
  trayectoria bajo condiciones explícitas;
\item
  caracterizar mediante VI alguna resolución suficiente para el estado
  futuro.
\end{enumerate}

Hasta aquí puede llegar la construcción sin especificar una dinámica que
aún no ha sido demostrada; queda abierta la pregunta que deberá resolverla:
\textbf{¿bajo qué instituciones una diferencia en la ejecución
intraperiódica sobrevive a la frontera como una diferencia en las
condiciones reproductivas del periodo siguiente?}


\cleardoublepage\section*{Conclusión general}\label{conclusiuxf3n-general}\markright{Conclusión general}\RBfronttitle{Conclusión general}\thispagestyle{plain}

El primer resultado es que el cierre terminal no agota la ejecución: una
misma red puede conservar sus saldos finales y producir necesidades
distributivas distintas cuando cambia el orden de los pagos. El segundo es
que esas necesidades individuales tampoco determinan por sí solas la
disponibilidad requerida, pues ésta depende de qué posiciones pueden
consolidarse y de cómo se sincronizan sus mínimos. A escala local, el efecto de
intercambiar dos obligaciones tampoco pertenece a la precedencia por sí sola:
una modificación de la arquitectura puede crear, eliminar o incluso invertir
la variación de \(M_0\) asociada al mismo intercambio. Este resultado ordena
requerimientos de referencia, no bienestar ni costes financieros realizados.
El tercero es que el requerimiento de referencia sólo se convierte en exposición efectiva después
de incorporar los recursos heredados. A partir de allí, la exposición fija un
piso operacional, pero el régimen contractual determina qué coberturas son
admisibles y la implementación financiera, qué posiciones persisten al
cierre; finalmente, una regla de frontera decide qué parte de ese estado
terminal se convierte en herencia.

La cadena I--V establece así relaciones de determinación condicionada, no una
causalidad automática: cada objeto resuelve una pregunta y deja abierta la que
exige la mediación siguiente. VI responde al problema general que aparece
cuando se intenta resumir alguno de esos objetos. Sean \(F\) y \(G\) dos
representaciones de una ejecución: \(F\) basta para determinar \(G\)
exactamente cuando nunca identifica dos casos que \(G\) necesita distinguir;
equivalentemente, \(G\) es constante sobre las fibras de \(F\) y factoriza a
través de ella. Las no-identificaciones obtenidas en los artículos anteriores
son instancias económicas de ese criterio.

Hasta aquí llega el resultado cerrado: el volumen determina cómo diferencias
de ejecución pueden producir diferencias de necesidad, exposición,
financiación, estado terminal y herencia monetario-financiera, pero todavía no
determina cuándo esas diferencias alteran el estado reproductivo real del
período siguiente. VII fija precisamente esa frontera: una dinámica conjunta
requiere construir el estado real-material y especificar su regla de
transición; después, demostrar bajo qué instituciones una diferencia
intraperiódica sobrevive al cierre como diferencia reproductiva en \(t+1\).

La construcción deja además una regla metodológica para su continuación: ningún
objeto económico nuevo debe introducirse antes de que una pregunta económica muestre
la insuficiencia del objeto disponible; sólo entonces corresponde elegir la representación
matemática que permita responderla.

El problema abierto puede formularse sin atribuir al análisis una conclusión
que todavía no establece: la secuenciación monetaria importa para la
reproducción interperiódica sólo cuando alguna de las diferencias que produce
atraviesa la frontera del período. Determinar las condiciones de esa
persistencia es el siguiente problema del programa.

\cleardoublepage
\section*{Bibliografía}\label{bibliografia-general}\markright{Bibliografía}\RBfronttitle{Bibliografía}\thispagestyle{plain}

\begin{refblock}
\begin{itemize}
\tightlist
\item
  Abdel-Wahab, H. M., y T. Kameda. ``Scheduling to Minimize Maximum
  Cumulative Cost Subject to Series-Parallel Precedence Constraints.''
  Operations Research 26(1) (1978): 141--158.
  \url{https://doi.org/10.1287/opre.26.1.141}.
\item
  Alexandrova Kabadjova, Biliana, et al.~``Intraday Liquidity around the
  World.'' BIS Working Papers No.~1089. Basel: Bank for International
  Settlements, April 2023.
\item
  Astuti, Valerio, Adriano Baldeschi, Luca Bastianelli, Giuseppe Bruno,
  Ajit Desai, Danica Marsden y Riccardo Russo. ``Liquidity Optimization
  in Gross Settlement Systems with Quantum Reordering: Application to
  TARGET2.'' Journal of Economic Dynamics and Control 190 (2026): 105385.
  \url{https://doi.org/10.1016/j.jedc.2026.105385}.
\item
  Basel Committee on Banking Supervision (BCBS). Monitoring Tools for
  Intraday Liquidity Management. Basel: Bank for International
  Settlements, April 2013.
\item
  Bech, Morten L., y Rod Garratt. ``The Intraday Liquidity Management
  Game.'' Journal of Economic Theory 109(2) (2003): 198--219.
  \url{https://doi.org/10.1016/S0022-0531(03)00016-4}.
\item
  Becher, Christopher, Marco Galbiati, y Merxe Tudela. ``The Timing and
  Funding of CHAPS Sterling Payments.'' Economic Policy Review 14(2)
  (2008): 113--133.
\item
  Blackwell, David. ``Comparison of Experiments.'' Proceedings of the
  Second Berkeley Symposium on Mathematical Statistics and Probability,
  1951: 93--102.
\item
  Campbell, Martha. ``Money in the Circulation of Capital.'' En
  Christopher J. Arthur y Geert Reuten (eds.), The Circulation of
  Capital: Essays on Volume Two of Marx's Capital, 129--157. London:
  Palgrave Macmillan, 1998.
\item
  Campbell, Martha. ``The Credit System.'' En Martha Campbell y Geert
  Reuten (eds.), The Culmination of Capital: Essays on Volume III of
  Marx's Capital, 212--227. Palgrave Macmillan, 2002.
\item
  Committee on Payment and Settlement Systems (CPSS). Real-Time Gross
  Settlement Systems. CPMI Papers No.~22. Basel: Bank for International
  Settlements, March 1997.
\item
  European Central Bank (ECB). Guideline of the European Central Bank of
  5 December 2012 on a Trans-European Automated Real-time Gross settlement
  Express Transfer system (TARGET2) (ECB/2012/27), 2012.
\item
  Freeman, Alan. ``Marx without Equilibrium.'' Capital \& Class 19(2)
  (1995): 49--89. \url{https://doi.org/10.1177/030981689505600103}.
\item
  Galbiati, Marco, y Kimmo Soramäki. ``Liquidity-saving mechanisms and
  bank behaviour.'' Bank of England Working Paper No.~400, 2010.
\item
  de Brunhoff, Suzanne. Marx on Money. Trad. Maurice J. Goldbloom; prefacio
  de Duncan K. Foley. New York: Urizen Books, 1976 {[}1967{]}.
\item
  Foley, Duncan K. ``The Value of Money, the Value of Labor Power and the
  Marxian Transformation Problem.'' Review of Radical Political Economics
  14(2) (1982a): 37--47. \url{https://doi.org/10.1177/048661348201400204}.
\item
  Foley, Duncan K. ``Realization and Accumulation in a Marxian Model of the
  Circuit of Capital.'' Journal of Economic Theory 28(2) (1982b): 300--319.
  \url{https://doi.org/10.1016/0022-0531(82)90063-1}.
\item
  Gehrke, Christian. ``Marx's Reproduction Schemes and Multi-Sector
  Growth Models.'' European Journal of the History of Economic Thought
  25(5) (2018): 859--892. \url{https://doi.org/10.1080/09672567.2018.1475500}.
\item
  Godley, Wynne and Marc Lavoie. Monetary Economics: An Integrated
  Approach to Credit, Money, Income, Production and Wealth. Palgrave
  Macmillan, 2007.
\item
  Graziani, Augusto. The Monetary Theory of Production. Cambridge
  University Press, 2003.
\item
  Hayakawa, Hitoshi. ``Does a Central Clearing Counterparty Reduce
  Liquidity Needs?'' Journal of Economic Interaction and Coordination
  13(1) (2018): 9--50.
  \url{https://doi.org/10.1007/s11403-017-0208-1}.
\item
  Hayakawa, Hitoshi. ``Liquidity in Financial Networks.'' Computational
  Economics 55(1) (2020): 253--301.
  \url{https://doi.org/10.1007/s10614-019-09895-x}.
\item
  Hayakawa, Hitoshi, Toshimasa Ishii, Hirotaka Ono y Yushi Uno.
  ``Settlement Fund Circulation Problem.'' Discrete Applied Mathematics
  265 (2019): 86--103; versión de conferencia: ISAAC 2017, LIPIcs 92,
  46:1--46:13. \url{https://doi.org/10.1016/j.dam.2019.03.017}.
\item
  Houston, Robin. ``Tackling the Minimal Superpermutation Problem.''
  arXiv:1408.5108, 2014.
\item
  Johnston, Nathaniel. ``Non-uniqueness of minimal superpermutations.''
  Discrete Mathematics 313(14) (2013): 1553--1557.
  \url{https://doi.org/10.1016/j.disc.2013.03.024}.
\item
  Kahn, Charles M., James McAndrews, y William Roberds. ``Settlement
  Risk under Gross and Net Settlement.'' Journal of Money, Credit and
  Banking 35(4) (2003): 591--608.
  \url{https://doi.org/10.1353/mcb.2003.0030}.
\item
  Kalecki, Michał. ``The Marxian Equations of Reproduction and Modern
  Economics.'' Social Science Information 7(6) (1968): 73--79.
  \url{https://doi.org/10.1177/053901846800700609}.
\item
  Kellerer, Hans, Vladimir Kotov, Franz Rendl y Gerhard J. Woeginger.
  ``The Stock Size Problem.'' Operations Research 46(3-supplement-3)
  (1998): S1--S12. \url{https://doi.org/10.1287/opre.46.3.S1}.
\item
  Kliman, Andrew, y Ted McGlone. ``A Temporal Single-System
  Interpretation of Marx's Value Theory.'' Review of Political Economy
  11(1) (1999): 33--59. \url{https://doi.org/10.1080/095382599107165}.
\item
  Lavoie, Marc. Post-Keynesian Economics: New Foundations. Edward Elgar,
  2014.
\item
  Leinster, Tom. Basic Category Theory. Cambridge: Cambridge University
  Press, 2014, cap. 3 ``Interlude on sets'', 65--82.
  \url{https://doi.org/10.1017/CBO9781107360068.005}.
\item
  Leontief, Wassily, y András Bródy. ``Money-flow Computations.''
  Economic Systems Research 5(3) (1993): 225--233.
  \url{https://doi.org/10.1080/09535319300000019}.
\item
  Luxemburg, Rosa. The Accumulation of Capital. London: Routledge, 2003
  {[}1913{]}.
\item
  Lowe, Adolph. ``A Bird's-Eye View of an Expanding Traverse.'' Cap. 12
  de The Path of Economic Growth, 109--114. Cambridge: Cambridge
  University Press, 1976.
  \url{https://doi.org/10.1017/CBO9780511898327.013}.
\item
  Lowe, Adolph. ``Some Comments on the Role of Working Capital in the
  Traverse.'' Cap. 17 de The Path of Economic Growth, 176--181.
  Cambridge: Cambridge University Press, 1976.
  \url{https://doi.org/10.1017/CBO9780511898327.018}.
\item
  Maldonado-Filho, Eduardo. ``Release and Tying up of Productive Capital
  and the `Transformation Problem'.'' Conference paper; versión
  presentada en la ASSA Conference, Washington, D.C., January 1995.
\item
  Marx, Karl. El capital. Crítica de la economía política. Libro
  primero: El proceso de producción del capital. Trad. Pedro Scaron.
  Siglo XXI, 1975 {[}1867{]}, cap. III, sec.~3(b), ``Medio de pago''.
\item
  Marx, Karl. El capital. Crítica de la economía política. Libro
  segundo: El proceso de circulación del capital. Trad. Pedro Scaron.
  Siglo XXI, 1976 {[}1885{]}, esp.~Parte III, caps. 18--21.
\item
  Marx, Karl. El capital. Crítica de la economía política. Libro
  tercero: El proceso global de la producción capitalista. Trad. León
  Mames; revisión y notas de Pedro Scaron. Siglo XXI, 1976 {[}1894{]},
  cap. 6, sec.~II, ``Apreciación, depreciación, liberación y vinculación
  de capital''.
\item
  Marx, Karl. Teorías sobre la plusvalía. Trad. Wenceslao Roces. México:
  Fondo de Cultura Económica, 1980 {[}manuscritos 1861--1863{]}, cap.
  XVII.
\item
  McAndrews, James J., y Samira Rajan. ``The Timing and Funding of
  Fedwire Funds Transfers.'' Economic Policy Review 6(2) (2000): 17--32.
\item
  McMahon, Christopher, Donald McGillivray, Ajit Desai, Francisco
  Rivadeneyra, Jean-Paul Lam, Thomas Lo, Danica Marsden y Vladimir
  Skavysh. ``Improving the Efficiency of Payments Systems Using Quantum
  Computing.'' Management Science 70(10) (2024): 7325--7341; versión de
  trabajo: Bank of Canada Staff Working Paper 2022-53, 2022.
\item
  Morishima, Michio, y Francis Seton. ``Aggregation in Leontief Matrices and
  the Labour Theory of Value.'' Econometrica 29(2) (1961): 203--220.
  \url{https://doi.org/10.2307/1909289}.
\item
  Morishima, Michio. Marx's Economics: A Dual Theory of Value and Growth.
  Cambridge: Cambridge University Press, 1973.
\item
  Minsky, Hyman P. ``The Financial Instability Hypothesis.'' Levy
  Economics Institute Working Paper No.~74, 1992.
\item
  Minsky, Hyman P. Stabilizing an Unstable Economy. Yale University
  Press, 1986.
\item
  Mitchell, William, L. Randall Wray and Martin Watts. Macroeconomics.
  Red Globe Press, 2019.
\item
  Monma, Clyde L. ``Sequencing to Minimize the Maximum Job Cost.''
  Operations Research 28(4) (1980): 942--951.
  \url{https://doi.org/10.1287/opre.28.4.942}.
\item
  Monma, Clyde L., y Jeffrey B. Sidney. ``Sequencing with Series-Parallel
  Precedence Constraints.'' Mathematics of Operations Research 4(3)
  (1979): 215--224. \url{https://doi.org/10.1287/moor.4.3.215}.
\item
  Pasinetti, Luigi L. Structural Change and Economic Growth: A
  Theoretical Essay on the Dynamics of the Wealth of Nations. Cambridge:
  Cambridge University Press, 1981.
\item
  Postnikov, Alexander. ``Permutohedra, Associahedra, and Beyond.''
  International Mathematics Research Notices 2009(6) (2009): 1026--1106.
  \url{https://doi.org/10.1093/imrn/rnn153}.
\item
  Räihä, Kari-Jouko and Esko Ukkonen. ``The shortest common
  supersequence problem over binary alphabet is NP-complete.''
  Theoretical Computer Science 16(2), 1981: 187--198.
\item
  Rotemberg, Julio J. ``Minimal Settlement Assets in Economies with
  Interconnected Financial Obligations.'' Journal of Money, Credit and
  Banking 43(1) (2011): 81--108; versión de trabajo: NBER Working Paper
  14222, 2008. \url{https://doi.org/10.1111/j.1538-4616.2010.00366.x}.
\item
  Stanley, Richard P. ``Acyclic Orientations of Graphs.'' Discrete
  Mathematics 5(2) (1973): 171--178.
  \url{https://doi.org/10.1016/0012-365X(73)90108-8}.
\item
  Tombazos, Stavros. Time in Marx: The Categories of Time in Marx's
  Capital. Brill, 2014 {[}orig. French 1994{]}.
\item
  Trigg, Andrew B. ``Reconstructing Marx's Theory of Credit and Payment
  Crises under Simple Circulation.'' Review of Political Economy 33(4)
  (2021): 746--768. \url{https://doi.org/10.1080/09538259.2021.1897751}.
\item
  Trigg, Andrew B. Marxian Reproduction Schema: Money and Aggregate
  Demand in a Capitalist Economy. Routledge, 2006.
\item
  Veronese Passarella, Marco. ``It Is Not la Vie en Rose: New Insights from
  Graziani's Theory of the Monetary Circuit.'' European Journal of Economics
  and Economic Policies: Intervention 21(3) (2024): 461--485.
  \url{https://doi.org/10.4337/ejeep.2022.0084}.
\item
  Veronese Passarella, Marco. ``Technical Change and the Monetary Circuit:
  An Input-Output Stock-Flow Consistent Dynamic Model.'' Department of
  Economics, University of Siena, Working Paper 903, 2023.
\item
  Wray, L. Randall. Modern Money Theory, 2nd ed.~Palgrave Macmillan,
  2015.
\item
  Ziegler, Günter M. Lectures on Polytopes. Graduate Texts in
  Mathematics 152. New York: Springer, 1995, esp.~``Fans, Arrangements,
  Zonotopes, and Tilings'', 191--230.
  \url{https://doi.org/10.1007/978-1-4613-8431-1}.
\end{itemize}
\end{refblock}

\cleardoublepage
\thispagestyle{empty}
\vspace*{\fill}
\begin{center}
{\small\scshape RB · Ensayos 2026 · 02\par}
\vspace{4mm}
{\footnotesize v1.01 · edición de circulación · octubre de 2026\par}
{\footnotesize Compuesto en formato A5 con EB Garamond y Latin Modern Math.\par}
{\footnotesize Ciudad de México.\par}
\end{center}
\vspace*{\fill}
\clearpage
\thispagestyle{empty}\null\clearpage
\end{document}
